% =====================================================================
% Capítulo 5. Desarrollo, análisis y resultados
% Estado: 5.1 redactada (06/10/2026); 5.2 redactada (15/10/2026, sesión R1);
% 5.3 redactada (04/11/2026, sesión R2); 5.4 y 5.5 redactadas (06/11/2026,
% sesión R3); 5.6 redactada (06/11/2026, sesión R4); 5.7 redactada (06/11/2026, sesión R4).
% Figuras de 5.1: definitivas (07/10/2026), generadas por src/14_exploracion.py
% directamente en memoria/figuras/.
% =====================================================================

\chapter{Desarrollo, análisis y resultados}\label{cap:resultados}

Este capítulo presenta los resultados del trabajo en el orden de los
objetivos específicos. Comienza con el análisis exploratorio de la serie de
demanda, que caracteriza sus propiedades estadísticas y fundamenta las
decisiones de modelización. Le sigue la validación en 2023, en la que se
elige la configuración de cada familia y se describe el comportamiento de los
cuatro modelos. A continuación se presentan la precisión puntual de los
cuatro modelos en el año de prueba, la calibración de sus intervalos de
predicción y el desglose de ambos por hora del día y tipo de día, y el
capítulo termina con la discusión de los resultados.

\section{Análisis exploratorio de la serie}\label{sec:exploratorio}

Esta sección responde al objetivo~OE2: caracterizar estadísticamente la
serie de demanda antes de modelizarla. Todo el análisis se realiza
exclusivamente con los datos de 2017-2022, el periodo de entrenamiento, que
comprende 52\,583 horas. Los años 2023 y 2024 no intervienen en ningún
cálculo, de modo que las decisiones que se derivan de esta sección no
incorporan información de los periodos de validación y prueba. Salvo en el
análisis de la relación con la temperatura, que se limita a los días
laborables fuera del periodo de la pandemia, los cálculos emplean la serie
completa de 2017-2022, porque las herramientas utilizadas requieren una serie
continua.

% Cifras generadas por src/14_exploracion.py (bloque 0, estadísticos
% descriptivos), con el dataset oficial (huella comprobada) y pandas 2.3.3.
En el periodo, la demanda tiene una media de 28\,049~MW y una desviación
típica de 4560~MW; excluido el periodo de la pandemia, la media sube a 28\,251~MW. El
valor horario mínimo, 16\,117~MW, y el máximo, 41\,829~MW, se analizan en la
sección~\ref{sec:explo-anomalias}. La tabla~\ref{tab:medias-anuales} muestra
la demanda media de cada año, que anticipa el comportamiento de la tendencia:
estable entre 2017 y 2019, con un descenso marcado en 2020, una recuperación
parcial en 2021 y un nuevo descenso en 2022, hasta el valor más bajo del
periodo.

\begin{table}[htbp]
	\centering
	\caption{Demanda media anual, 2017-2022}
	\label{tab:medias-anuales}
	\small
	\begin{tabular}{lrrrrrr}
		\hline
		\textbf{Año} & 2017 & 2018 & 2019 & 2020 & 2021 & 2022 \\
		\hline
		\textbf{Demanda media (MW)} & 28\,849 & 29\,056 & 28\,527 & 27\,078 & 27\,838 & 26\,946 \\
		\hline
	\end{tabular}
\end{table}

El resto de la sección sigue el orden del análisis. Primero se identifican las
periodicidades de la serie mediante el periodograma y se separan sus
componentes con la descomposición MSTL. Después se estudian la
estacionariedad, la autocorrelación, la relación con la temperatura y los
efectos del calendario y de los episodios atípicos. La sección termina con una
síntesis de las implicaciones de estos hallazgos para la modelización.

% ---------------------------------------------------------------------
\subsection{Periodicidades: el periodograma}\label{sec:explo-periodograma}

El primer paso es identificar qué ciclos contiene la serie y cuánto pesa cada
uno. La figura~\ref{fig:periodograma} muestra el periodograma de la demanda
horaria, que reparte la varianza de la serie entre ciclos de distinta
duración (sección~\ref{sec:espectral}). El eje horizontal indica la duración
del ciclo, en horas, y el vertical la potencia espectral, que es proporcional al
cuadrado de la amplitud del ciclo y se mide en MW$^{2}\cdot$h; ambos ejes
están en escala logarítmica.
Las líneas discontinuas señalan los periodos de 24, 168 y 8766 horas, es
decir, el día, la semana y el año medio de 365,25 días.

\begin{figure}[htbp]
	\centering
	\includegraphics[width=\textwidth]{periodograma.png}
	\caption{Periodograma de la demanda horaria, 2017-2022. Ambos ejes en
	escala logarítmica; las líneas discontinuas marcan los periodos de 24, 168
	y 8766 horas.}
	\label{fig:periodograma}
\end{figure}

Las dos frecuencias más potentes de toda la serie son exactamente las de 24 y
168 horas. Les siguen picos en 12, 84 y 28 horas, que no son ciclos nuevos,
sino armónicos de los anteriores: 12 es la mitad de 24, y 84 y 28 son la
mitad y la sexta parte de 168. Su presencia indica que los ciclos diario y
semanal no tienen forma sinusoidal, como cabe esperar de un perfil diario con
un valle nocturno largo y dos máximos, y de una semana en la que el sábado y
el domingo se separan bruscamente de los laborables.

En los periodos largos destacan las frecuencias de 4381,9 horas, el
semestre, y de 1460,6 horas, la sexta parte del año, que es un armónico
anual. El periodo anual propiamente dicho, que en la rejilla de frecuencias
del periodograma corresponde a 8763,8 horas, ocupa el décimo puesto de 26\,291
frecuencias, pero su potencia es solo el 9,2\,\% de la del semestre, lo que
equivale a una amplitud de aproximadamente un 30\,\% de la semestral. La
explicación está en la forma del ciclo anual: la demanda tiene dos máximos al
año, uno en invierno y otro en verano, y dos mínimos en primavera y otoño, de
modo que la oscilación dominante tiene periodo semestral.

En síntesis, el periodograma muestra picos en los tres periodos
estacionales: 24 y 168 horas son las dos frecuencias más potentes de la
serie, y el periodo anual ocupa el décimo puesto, aunque su potencia es
inferior a la de su segundo armónico, de 4382 horas, reflejo del doble
máximo invernal y estival. Esta observación tiene una consecuencia directa
para la modelización: si la estacionalidad anual se representa con términos
de Fourier, un solo armónico no basta, y hacen falta al menos dos
(sección~\ref{sec:fourier}).

El análisis tiene dos limitaciones. La primera es la resolución en los
periodos largos. Con 52\,583 horas, las frecuencias contiguas a la anual
corresponden a periodos de 7512 y 10\,517 horas, de modo que el periodograma
no puede distinguir ciclos cercanos al año y la estimación de su potencia es
poco precisa. La segunda es la fuga espectral desde la tendencia: los cambios
de nivel entre años reparten potencia hacia las frecuencias más bajas, lo que
explica la subida de la curva en el extremo derecho de la figura.

% ---------------------------------------------------------------------
\subsection{Descomposición en tendencia y estacionalidades}\label{sec:explo-mstl}

Una vez identificados los periodos, la descomposición MSTL
(sección~\ref{sec:mstl}) separa la serie en una tendencia, tres componentes
estacionales de 24, 168 y 8766 horas y un residuo, cuya suma reproduce la
serie original. Se emplea la implementación de la librería statsmodels con
sus parámetros por defecto. La figura~\ref{fig:mstl} muestra la serie y sus
componentes, todas en megavatios y con el tiempo en el eje horizontal, y la
tabla~\ref{tab:mstl} resume la variabilidad de cada componente mediante su
desviación típica.

\begin{figure}[htbp]
	\centering
	\includegraphics[width=\textwidth]{descomposicion_mstl.png}
	\caption{Descomposición MSTL de la demanda horaria, 2017-2022. De arriba
	abajo: serie observada, tendencia, componentes estacionales diaria (24~h),
	semanal (168~h) y anual (8766~h), y residuo. Todas las componentes en MW.}
	\label{fig:mstl}
\end{figure}

\begin{table}[htbp]
	\centering
	\caption{Desviación típica de las componentes de la descomposición MSTL, 2017-2022}
	\label{tab:mstl}
	\small
	\begin{tabular}{lr}
		\hline
		\textbf{Componente} & \textbf{Desviación típica (MW)} \\
		\hline
		Estacional diaria (24~h) & 3336 \\
		Estacional semanal (168~h) & 2070 \\
		Estacional anual (8766~h) & 1816 \\
		Residuo & 1052 \\
		Tendencia & 828 \\
		\hline
	\end{tabular}
\end{table}

El ciclo diario es la mayor fuente de variación de la serie, seguido del
semanal y del anual, que tienen un peso parecido entre sí. El residuo, que
recoge todo lo que no explican la tendencia y los ciclos regulares, tiene
una desviación típica de 1052~MW, en torno al 4\,\% de la demanda media. Es
en ese residuo donde actúan la temperatura, que se aparta de su valor
estacional típico, y el calendario, que altera la regularidad semanal.

La tendencia confirma lo que anticipaban las medias anuales: un nivel
estable entre 2017 y 2019, en torno a 29\,000~MW, una caída en 2020, una
recuperación parcial en 2021 y un descenso fuerte a lo largo de 2022. El
tramo final debe interpretarse con cautela, porque el suavizado en el que se
basa la descomposición es menos fiable en los extremos de la serie.

La componente anual presenta una doble joroba: un máximo en enero, uno
secundario en julio y mínimos en primavera y otoño, coherente con el
periodograma. Sobre ella aparecen agujas descendentes en las mismas fechas de
cada año, que corresponden a festivos de fecha fija como Navidad o Año Nuevo.
Como la descomposición no conoce el calendario, reparte su efecto entre la
componente anual, cuando el festivo cae siempre en la misma fecha, y el
residuo.

Las componentes diaria y semanal no tienen amplitud constante. La diaria es
mayor en invierno y en verano que en primavera y otoño, es decir, el perfil
diario cambia con la estación, y la semanal se reduce de forma visible en la
primavera de 2020, cuando las restricciones de la pandemia redujeron la
diferencia entre los laborables y el fin de semana. Esta variación tiene
consecuencias para la modelización: los términos de Fourier imponen un
patrón que se repite sin cambios (sección~\ref{sec:fourier}), mientras que
la descomposición muestra que el patrón real cambia a lo largo del año.

La principal limitación de la descomposición es que, con solo seis ciclos
anuales, la componente anual se estima con poca información y debe
considerarse aproximada. La descomposición es, además, una herramienta
descriptiva: no forma parte de ningún modelo de previsión del trabajo.

% ---------------------------------------------------------------------
\subsection{Estacionariedad}\label{sec:explo-estacionariedad}

Los contrastes de estacionariedad determinan si la serie debe diferenciarse
antes de modelizarla. Se aplican los contrastes ADF y KPSS
(sección~\ref{sec:raiz-unitaria}), que tienen hipótesis nulas opuestas, a la
serie de demanda de 2017-2022. En ambos se incluye una constante, sin
tendencia lineal. La tabla~\ref{tab:estacionariedad} recoge los resultados.

\begin{table}[htbp]
	\centering
	\caption{Contrastes de estacionariedad sobre la demanda horaria, 2017-2022}
	\label{tab:estacionariedad}
	\small
	\begin{tabularx}{\textwidth}{>{\raggedright\arraybackslash}p{4.2cm} >{\raggedright\arraybackslash}X r r}
		\hline
		\textbf{Contraste} & \textbf{Hipótesis nula} & \textbf{Estadístico} & \textbf{$p$-valor} \\
		\hline
		ADF, 200 retardos elegidos por AIC & Raíz unitaria & $-6{,}96$ & $\approx 0$ \\
		ADF, 336 retardos fijos & Raíz unitaria & $-5{,}50$ & $\approx 0$ \\
		KPSS, constante & Estacionariedad en nivel & 6,21 & $< 0{,}01$ \\
		\hline
	\end{tabularx}
\end{table}

El contraste ADF rechaza con claridad la hipótesis de raíz unitaria. El
número de retardos se eligió por el criterio de Akaike con un máximo de 200,
y el criterio eligió precisamente ese máximo, lo que indica que podría
preferir más. Para comprobar que el resultado no depende de ello, se repitió
el contraste con 336 retardos fijos, dos semanas completas, y el rechazo se
mantiene. El contraste KPSS, por su parte, rechaza la estacionariedad en
nivel: su estadístico, 6,21, multiplica por más de ocho el valor crítico al
1\,\%, que es 0,739. El $p$-valor exacto es inferior a 0,01, el menor valor
que tabula la implementación empleada.

Que ambos contrastes rechacen su hipótesis nula no es una contradicción. La
serie no tiene una raíz unitaria, es decir, sus perturbaciones no se
acumulan indefinidamente, pero tampoco oscila en torno a una media constante:
su media cambia con la hora, el día de la semana, la estación y la tendencia.
El criterio de decisión se fijó antes de realizar los contrastes: en este
caso, en lugar de diferenciar la serie, se modelan explícitamente la
tendencia, las estacionalidades y el efecto de las variables exógenas, y la
estacionariedad se vuelve a contrastar sobre los residuos del SARIMAX
(objetivo~OE3; sección~\ref{sec:val-sarimax}). Como el protocolo de
comparación no admite variables de tendencia o de año, la tendencia se recoge
mediante una variable de nivel construida con la demanda de los días
anteriores a cada previsión, que se describe en la
sección~\ref{sec:esp-sarimax}.

Los resultados tienen dos matices. Con más de 52\,000 observaciones, los
contrastes tienen una potencia tan alta que detectan desviaciones mínimas de
la hipótesis nula, de modo que un rechazo no implica por sí solo una
desviación de importancia práctica. Y no se han contrastado las raíces
unitarias estacionales (sección~\ref{sec:raiz-unitaria}): los contrastes
aplicados se refieren a la raíz unitaria ordinaria, y no permiten distinguir
si la estacionalidad es determinista o estocástica. Esta limitación se
retoma en las conclusiones.

% ---------------------------------------------------------------------
\subsection{Autocorrelación}\label{sec:explo-acf}

La figura~\ref{fig:acf} muestra las funciones de autocorrelación simple
(ACF) y parcial (PACF) de la demanda y del residuo de la descomposición MSTL,
hasta el retardo de 400 horas. En cada panel, el eje horizontal es el
retardo, en horas, y el vertical el coeficiente de correlación, entre $-1$ y
1. La banda sombreada delimita, al 95\,\%, los valores compatibles con una
correlación nula en cada retardo. En los paneles de la PACF es la banda de
ruido blanco de la sección~\ref{sec:acf}, de $\pm 0{,}0085$, tan estrecha que
apenas se distingue. En los de la ACF, la implementación empleada aplica la
fórmula de Bartlett, que acumula la incertidumbre de las correlaciones de los
retardos anteriores, por lo que la banda se ensancha con el retardo hasta unos
$\pm 0{,}08$ para la demanda y $\pm 0{,}10$ para el residuo a 400 horas.
Las líneas discontinuas marcan los retardos de 24, 168 y 336 horas.

\begin{figure}[htbp]
	\centering
	\includegraphics[width=\textwidth]{acf_pacf.png}
	\caption{Autocorrelación simple (izquierda) y parcial (derecha) de la
	demanda (arriba) y del residuo de la descomposición MSTL (abajo),
	2017-2022, hasta 400 horas de retardo. Líneas discontinuas en 24, 168 y
	336 horas.}
	\label{fig:acf}
\end{figure}

La ACF de la demanda oscila con periodo de 24 horas y su envolvente alcanza
los máximos en los múltiplos de 168. La tabla~\ref{tab:acf} recoge los
valores en los retardos más relevantes. La correlación con la misma hora de
la semana anterior (0,90) es mayor que con la misma hora del día anterior
(0,83): la demanda de un martes a las 10:00 se parece más a la del martes
anterior a esa hora que a la del lunes. Este resultado respalda el uso del
ingenuo estacional semanal como modelo de referencia y del retardo de 168
horas como uno de los más informativos, y tiene especial relevancia porque
el retardo de 24 horas no está disponible para todas las horas del día
previsto, al emitirse la previsión a las 12:00 del día anterior.

\begin{table}[htbp]
	\centering
	\caption{Autocorrelación simple en retardos seleccionados, 2017-2022}
	\label{tab:acf}
	\small
	\begin{tabular}{lrrrr}
		\hline
		\textbf{Serie} & \textbf{1~h} & \textbf{24~h} & \textbf{168~h} & \textbf{336~h} \\
		\hline
		Demanda & 0,96 & 0,83 & 0,90 & 0,84 \\
		Residuo de MSTL & 0,99 & 0,75 & 0,31 & 0,03 \\
		\hline
	\end{tabular}
\end{table}

La PACF de la demanda, que mide la relación directa con cada retardo una vez
descontados los intermedios, concentra los valores relevantes en tres
grupos: los retardos de 1 y 2 horas, los de 24 a 26 y los de 168 y 169. La
serie depende directamente, por tanto, de su pasado inmediato, del mismo
momento del día anterior y del mismo momento de la semana anterior, lo que
coincide con la estructura de estacionalidad múltiple identificada en el
periodograma.

El residuo de la descomposición se comporta de otra forma. Sus picos
estacionales han desaparecido, señal de que la descomposición ha extraído
los ciclos regulares, pero su autocorrelación es muy persistente: 0,99 a una
hora, 0,75 a un día y todavía 0,31 a una semana, con un descenso lento y sin
oscilaciones. Su PACF, con dos valores destacados en los retardos 1 y 2 y
casi nada después, corresponde a un proceso autorregresivo de orden dos. La
interpretación es que, una vez descontados los ciclos, la demanda se aparta
de su patrón típico en episodios que duran varios días, como una ola de
frío o de calor o un periodo festivo. La hipótesis que se deriva para el
objetivo~OE3, todavía sin comprobar, es que la temperatura y el calendario
absorberán buena parte de esa persistencia, y que lo que quede podrá
recogerlo un proceso ARMA de orden bajo en los errores del SARIMAX.

% ---------------------------------------------------------------------
\subsection{Relación con la temperatura}\label{sec:explo-temperatura}

La relación entre la demanda y la temperatura se analiza a escala diaria y
solo con los días laborables de 2017-2022 fuera del periodo de la pandemia,
1410 días, para que el efecto del calendario no se confunda con el de la
temperatura. La figura~\ref{fig:temperatura} representa cada día como un
punto, con la temperatura media diaria ponderada por población en el eje
horizontal, en grados Celsius, y la demanda media diaria en el vertical, en
megavatios. La línea une las medias de la demanda en intervalos de 1~°C.

\begin{figure}[htbp]
	\centering
	\includegraphics[width=0.85\textwidth]{demanda_temperatura.png}
	\caption{Demanda media diaria frente a temperatura media diaria ponderada
	por población en los días laborables de 2017-2022, excluido el periodo de
	la pandemia. La línea une las medias por intervalos de 1~°C.}
	\label{fig:temperatura}
\end{figure}

La relación tiene forma de U, coherente con la respuesta no lineal de la
demanda española a la temperatura descrita en la literatura
\parencite{moralcarcedo2005modelling}. La demanda media alcanza su mínimo en
torno a 16~°C, con 27\,299~MW, y se mantiene casi plana entre 15 y 20~°C. Por
debajo y por encima de esa zona crece a un ritmo parecido en ambas ramas,
del orden de 550 a 600~MW por grado. La dispersión vertical de los puntos
para una misma temperatura es amplia, de varios miles de megavatios, y
recoge el efecto de la tendencia entre años, que la figura no descuenta, y de
otros factores.

La figura~\ref{fig:ajuste-gd} muestra el resultado del procedimiento de
selección de umbrales descrito en la sección~\ref{sec:datos-gradosdia}. El
óptimo se alcanza en $T_c = 15{,}5$~°C para la calefacción y
$T_r = 19{,}5$~°C para la refrigeración, en el interior de la rejilla
explorada, de modo que se cumple el criterio de validez fijado de antemano.
La superficie de error es, sin embargo, bastante plana cerca del óptimo: las
parejas con $T_c$ entre 15 y 16,5~°C y $T_r$ entre 18,5 y 20~°C difieren del
mínimo en menos de un 0,35\,\% de error, y ese intervalo es el que debe
considerarse como resultado, más que la pareja exacta. Con los umbrales
elegidos, cada grado-día ponderado de calefacción aumenta la demanda media
diaria en unos 610~MW y cada grado-día de refrigeración en unos 596~MW. El
modelo explica el 71,7\,\% de la varianza de la demanda diaria de los
laborables ($R^2 = 0{,}717$), cifra que incluye la contribución de las
variables indicadoras de año.

\begin{figure}[htbp]
	\centering
	\includegraphics[width=0.85\textwidth]{demanda_temperatura_ajuste.png}
	\caption{Demanda media diaria observada y ajustada con grados-día de
	calefacción y de refrigeración y una variable indicadora por año, en los
	días laborables de 2017-2022 sin el periodo de la pandemia. Las líneas
	discontinuas marcan los umbrales $T_c = 15{,}5$~°C y $T_r = 19{,}5$~°C.}
	\label{fig:ajuste-gd}
\end{figure}

La figura permite tres observaciones. Los valores ajustados forman bandas
paralelas, una por año, porque cada año tiene su propio nivel. Entre los dos
umbrales, el ajuste no es plano, como cabría esperar si la temperatura
peninsular determinara por sí sola los grados-día: lo que ocurre es que,
aunque la temperatura media ponderada esté en la zona neutra, algunas
provincias están fuera de ella. Es la misma situación del ejemplo de la
sección~\ref{sec:datos-gradosdia} y confirma la utilidad de calcular los
grados-día por provincia. Por último, en la rama fría, entre 8 y 12~°C,
numerosos laborables quedan por encima del ajuste. Una posible explicación
es el mayor uso de iluminación en los días cortos del invierno, y otra, que
la respuesta al frío no sea exactamente lineal; ninguna de las dos se ha
comprobado.

% ---------------------------------------------------------------------
\subsection{Calendario y episodios atípicos}\label{sec:explo-anomalias}

El ajuste anterior permite localizar los días laborables cuya demanda queda
más lejos de lo que predice la temperatura. Los veinte con mayor déficit
presentan una demanda entre 4000 y 6300~MW inferior a la ajustada, y todos
tienen una explicación de calendario: dieciséis pertenecen al periodo
navideño, entre el 23 de diciembre y el 2 de enero, y los cuatro restantes
son lunes de Pascua. Ninguno es de agosto, pese a ser el mes de vacaciones
por excelencia, lo que sugiere que su efecto ya queda recogido por la
temperatura y por el nivel anual.

Los dos grupos son diferentes desde el punto de vista de la modelización. El
lunes de Pascua es festivo en varias comunidades, de modo que la fracción de
población en festivo recoge al menos parte de su efecto; la regresión de la
sección anterior no incluía esa variable. El periodo navideño, en cambio, no
lo recoge ninguna variable del conjunto de datos: sus días laborables son
días de trabajo a efectos legales, pero con una actividad muy inferior a la
habitual. Esta evidencia, obtenida solo con 2017-2022, justifica la variable
de periodo navideño que el protocolo de comparación añade a los modelos.

La borrasca Filomena, en enero de 2021, constituye el episodio meteorológico
más extremo del periodo. Durante doce días, la temperatura media diaria
ponderada se mantuvo entre 3,3 y 4,8~°C, y la demanda media de los días
laborables alcanzó entre 34\,000 y 35\,200~MW. El máximo horario de todo el
periodo, 41\,829~MW, se registró el viernes 8 de enero de 2021 a las 13:00.

El mínimo horario, 16\,117~MW, se registró el domingo 25 de diciembre de 2022
a las 05:00: una madrugada de domingo, de Navidad y de un año de demanda
baja. Es, además, uno de los festivos en domingo que el calendario no marca
(sección~\ref{sec:datos-calendario}). La comprobación de esos festivos
confirma que las ocho fiestas nacionales de fecha fija que cayeron en
domingo entre 2017 y 2022 tienen una fracción de población en festivo nula
en el conjunto de datos.

Ninguno de los valores revisados resulta erróneo. Todos los episodios
atípicos tienen una causa identificable y real, por lo que se mantienen en el
conjunto de datos sin corrección (sección~\ref{sec:datos-anomalias}).

% ---------------------------------------------------------------------
\subsection{Síntesis: implicaciones para la modelización}\label{sec:explo-sintesis}

La tabla~\ref{tab:sintesis-oe2} resume los hallazgos del análisis
exploratorio y la decisión de modelización que se deriva de cada uno. Todas
las decisiones se apoyan exclusivamente en datos de entrenamiento.

\begin{table}[htbp]
	\centering
	\caption{Hallazgos del análisis exploratorio y sus implicaciones para la modelización}
	\label{tab:sintesis-oe2}
	\small
	\begin{tabularx}{\textwidth}{>{\raggedright\arraybackslash}X >{\raggedright\arraybackslash}X}
		\hline
		\textbf{Hallazgo} & \textbf{Implicación} \\
		\hline
		Picos en 24 y 168~h; la potencia anual es inferior a la semestral & Términos de Fourier para los tres periodos, con al menos dos armónicos anuales; el número de cada periodo se elige por criterio de información (sección~\ref{sec:esp-sarimax}) \\
		La amplitud de los ciclos diario y semanal cambia con la estación & Los términos de Fourier fijos son una limitación del SARIMAX, que se atenúa con interacciones entre la temperatura y la hora (sección~\ref{sec:esp-sarimax}); los modelos de aprendizaje pueden capturarlo \\
		ADF rechaza la raíz unitaria y KPSS la estacionariedad en nivel & No se diferencia la serie: se modelan tendencia (mediante una variable de nivel construida con la demanda pasada), estacionalidad y exógenas, y se contrastan los residuos \\
		La ACF es mayor en 168~h que en 24~h & Ingenuo estacional semanal como referencia; el retardo de 168~h, entre los principales disponibles \\
		Residuo de MSTL persistente, de tipo AR(2) & Episodios de varios días que deben explicar la temperatura y el calendario; dinámica ARMA en los errores \\
		Relación en U con la temperatura & Grados-día por provincia con $T_c = 15{,}5$~°C y $T_r = 19{,}5$~°C como exógenas comunes \\
		Mayores déficits en el periodo navideño y el lunes de Pascua & Variable de periodo navideño; fracción de población en festivo \\
		Extremos con causa real (Filomena, Navidad) & No se corrigen; se vigilan en el desglose de resultados \\
		\hline
	\end{tabularx}
\end{table}

Con ello se cumple el criterio establecido para el objetivo~OE2 en la
tabla~\ref{tab:criterios}: la memoria presenta la descomposición de la serie,
el periodograma con los picos en los tres periodos estacionales, los
resultados de los contrastes ADF y KPSS y la curva demanda-temperatura
ajustada con grados-día.

% =====================================================================
% Sección 5.2 redactada el 15/10/2026 (sesión R1). Cifras de
% src/18_sarimax_validacion.py, src/20_lightgbm_validacion.py,
% src/23_nhits_validacion.py y src/24_validacion_desgloses.py
% (protocolo, apartados 10.10, 11.4, 12.6 y 13.8).
% =====================================================================
\section{Validación en 2023}\label{sec:validacion}

Esta sección presenta los resultados del año de validación, 2023: la
selección de la configuración de cada familia, el diagnóstico de los residuos
del SARIMAX que exige el objetivo~OE3 y la comparación de los cuatro modelos
seleccionados. Conviene recordar el papel de estas cifras. Sirven para elegir
la configuración de cada familia y para describir cómo se comporta cada
modelo, pero no son el rendimiento que se informa: las tres configuraciones
ganadoras se han elegido precisamente por su error en 2023, que es por ello
una estimación optimista (sección~\ref{sec:prot-particiones}). El rendimiento
de los modelos es el de 2024, que se presenta más adelante en este capítulo.

La sección comienza con la selección dentro de cada familia y el diagnóstico
del SARIMAX. Después compara los cuatro modelos en el conjunto del año,
analiza cómo se reparte el error a lo largo del día y del año y en los días
más difíciles, y termina con una síntesis.

% ---------------------------------------------------------------------
\subsection{Selección en cada familia}\label{sec:val-seleccion}

La tabla~\ref{tab:seleccion} resume la búsqueda de cada familia. LightGBM y
N-HiTS agotaron su presupuesto; el SARIMAX evaluó 17 de las 20
configuraciones posibles, por el diseño de su rejilla
(sección~\ref{sec:esp-sarimax}). Ninguna configuración falló.

\begin{table}[htbp]
	\centering
	\caption{Resumen de la selección en 2023 (MAE y sesgo en MW)}
	\label{tab:seleccion}
	\small
	\begin{tabularx}{\textwidth}{l *{5}{>{\raggedleft\arraybackslash}X}}
		\hline
		\textbf{Familia} & \textbf{Config.} & \textbf{MAE mínimo} & \textbf{MAE máximo} & \textbf{Sesgo de la ganadora} & \textbf{Tiempo de la búsqueda} \\
		\hline
		SARIMAX & 17 & 943,4 & 1888,2 & $-346$ & $\approx 1$~h \\
		LightGBM & 20 & 585,2 & 656,3 & $-142$ & unos 5~min \\
		N-HiTS & 20 & 716,4 & 936,6 & $+188$ & 82~min \\
		\hline
	\end{tabularx}
\end{table}

\paragraph{SARIMAX.} La tabla~\ref{tab:sarimax-rejilla} recoge las 17
configuraciones evaluadas. Su lectura más importante está en las tres
primeras filas, las de la configuración de referencia ARMA(2,1) antes y
después de cada enmienda: sin la variable de nivel, el modelo falla más que el
ingenuo (1888~MW) y sobrestima la demanda en 1660~MW de media; con ella, el
MAE baja a 1315~MW; y con los retardos de 48 y 168 horas, a 974~MW. Lo
decisivo para el SARIMAX fueron, por tanto, el nivel y los retardos, más que
el orden del proceso de los errores. Con ambos, las seis mejores
configuraciones quedan entre 943 y 954~MW, a menos de un 1,2\,\% de distancia,
de modo que el orden del ARMA es un ajuste fino. La ganadora es el ARMA(5,1),
con 943,4~MW. Cinco configuraciones forman un segundo grupo, con errores de
1059 a 1112~MW y sesgos de entre $-111$ y $+179$~MW, frente a los de entre
$-290$ y $-390$~MW de las demás; una hipótesis no comprobada es que el
procedimiento de estimación llega en ellas a una solución distinta. La ganadora por RMSE es también el ARMA(5,1), lo que indica
que la elección no depende de la métrica.

\begin{table}[htbp]
	\centering
	\caption{Configuraciones del SARIMAX evaluadas en 2023 (MW). Las tres primeras filas recogen la configuración de referencia ARMA(2,1) en cada fase; la tercera forma parte también de la rejilla final, cuyas quince configuraciones se ordenan después por MAE}
	\label{tab:sarimax-rejilla}
	\small
	\begin{tabular}{llrrr}
		\hline
		\textbf{Errores} & \textbf{Variables de la demanda} & \textbf{MAE} & \textbf{RMSE} & \textbf{Sesgo} \\
		\hline
		ARMA(2,1) & Ninguna (referencia inicial) & 1888,2 & 2454,5 & $-1660$ \\
		ARMA(2,1) & Nivel & 1314,8 & 1785,7 & $-629$ \\
		ARMA(2,1) & Nivel y retardos & 973,7 & 1471,4 & $-350$ \\
		\hline
		\textbf{ARMA(5,1)} & \textbf{Nivel y retardos} & \textbf{943,4} & \textbf{1414,4} & $\mathbf{-346}$ \\
		ARMA(4,2) & Nivel y retardos & 947,1 & 1416,8 & $-351$ \\
		ARMA(1,1) & Nivel y retardos & 951,9 & 1435,3 & $-291$ \\
		ARMA(4,0) & Nivel y retardos & 952,8 & 1438,3 & $-315$ \\
		ARMA(3,1) & Nivel y retardos & 953,2 & 1439,0 & $-316$ \\
		ARMA(3,0) & Nivel y retardos & 954,0 & 1440,6 & $-319$ \\
		ARMA(2,2) & Nivel y retardos & 965,5 & 1459,8 & $-340$ \\
		ARMA(1,2) & Nivel y retardos & 969,0 & 1464,4 & $-346$ \\
		ARMA(2,0) & Nivel y retardos & 1012,5 & 1524,3 & $-387$ \\
		ARMA(5,2) & Nivel y retardos & 1059,3 & 1575,7 & $-111$ \\
		ARMA(5,0) & Nivel y retardos & 1093,5 & 1641,2 & $+122$ \\
		ARMA(3,2) & Nivel y retardos & 1100,9 & 1644,1 & $+122$ \\
		ARMA(4,1) & Nivel y retardos & 1106,4 & 1655,7 & $+170$ \\
		ARMA(1,0) & Nivel y retardos & 1111,5 & 1671,1 & $+179$ \\
		\hline
	\end{tabular}
\end{table}

\paragraph{LightGBM.} Las 20 configuraciones quedan entre 585,2 y 656,3~MW,
de modo que la peor falla un 12\,\% más que la mejor, frente al 18\,\% de la
rejilla final del SARIMAX, y las siete mejores quedan entre 585 y 596~MW, a
menos de un 2\,\%. Ninguna llegó al máximo de árboles. Lo que más
pesa es la tasa de aprendizaje: las cinco mejores configuraciones tienen tasas
de entre 0,029 y 0,032, y las tres peores, de entre 0,15 y 0,19. La función de
pérdida, cuadrática o absoluta, no muestra una diferencia clara. Todas las
configuraciones tienen sesgo negativo, entre $-93$ y $-179$~MW. La ganadora
obtiene un MAE de 585,2~MW. Por ganancia en los árboles, las variables más
importantes son el retardo de 168 horas, con el 58,7\,\% del total, y el de
336 horas, con el 11,5\,\%. Esta medida debe leerse con cuidado: la
temperatura suma solo un 3\,\%, pero eso no significa que su efecto sea
pequeño, porque la demanda retardada ya contiene el efecto de la temperatura
de entonces, y el modelo solo necesita la temperatura para corregir lo que
cambia entre aquel día y el previsto.

\paragraph{N-HiTS.} Las 20 configuraciones quedan entre 716,4 y 936,6~MW, y
las ocho mejores, entre 716 y 767~MW. Doce llegan al máximo de 1000 pasos,
incluida la ganadora. La observación hecha al sortear el espacio, que las
tasas de aprendizaje muy bajas darían redes poco entrenadas, se cumple solo
en parte: de las ocho configuraciones con tasas menores que 0,0002, cuatro
quedan en la mitad inferior de la clasificación, pero otras dos quedan
segunda y quinta. Todas las
configuraciones tienen sesgo positivo, entre $+91$ y $+562$~MW, al contrario
que las del SARIMAX y LightGBM: N-HiTS se queda corto de forma sistemática.

La ganadora, reentrenada con las semillas 43 y 44, obtiene un MAE de 771,1 y
704,6~MW, frente a 716,4~MW con la semilla 42: un rango de 66,4~MW, con una
media de 730,7~MW. La distancia entre la ganadora y la segunda configuración,
de 723,4~MW, es de solo 7~MW, mucho menor que ese rango. El orden entre las
mejores configuraciones de N-HiTS está, por tanto, dentro de la variabilidad
debida a la semilla, y una diferencia de unas decenas de megavatios entre dos
configuraciones de la red no debe interpretarse como una diferencia real. La
previsión oficial sigue siendo la de la semilla 42
(sección~\ref{sec:prot-presupuesto}).

% ---------------------------------------------------------------------
\subsection{Diagnóstico del SARIMAX}\label{sec:val-sarimax}

El objetivo~OE3 exige contrastar si los residuos del SARIMAX se comportan
como ruido blanco (tabla~\ref{tab:criterios}). Sobre los residuos
estandarizados del ajuste de 2017-2022 de la configuración ganadora, el
contraste de Ljung-Box rechaza la ausencia de autocorrelación hasta los
retardos de 24 y de 168 horas: el estadístico vale $Q = 9061$ con 24
retardos y $Q = 17\,850$ con 168, con $p$-valores inferiores a 0,0001 en
ambos casos. Los grados de libertad empleados son 24 y 168; descontar los
seis parámetros del proceso ARMA (sección~\ref{sec:box-jenkins}) no cambia
la conclusión. Con unas 49\,000 horas de residuos, el contraste rechaza ante
autocorrelaciones muy pequeñas, de modo que lo informativo es su magnitud. La
autocorrelación a 168 horas es prácticamente nula, $-0{,}04$, pero a 24 horas
es de 0,36: el modelo aún tiende a repetir el error de la misma hora del día
anterior. El contraste ADF
rechaza con claridad la raíz unitaria, con un estadístico de $-36{,}6$, y el
KPSS rechaza la estacionariedad en nivel, con 2,9 y un $p$-valor inferior a
0,01, lo que indica que en los residuos quedan cambios de nivel lentos.

El resultado contrasta con la hipótesis formulada en el análisis
exploratorio (sección~\ref{sec:explo-acf}): que la temperatura y el
calendario absorberían buena parte de la persistencia del residuo de la
descomposición MSTL y que un proceso ARMA de orden bajo bastaría para el
resto. No fue así. Hicieron falta la variable de nivel y los retardos de 48 y
168 horas, el proceso de los errores ganador es de orden cinco, y aun así
queda una dependencia diaria que la especificación no recoge. Una causa
posible, no comprobada, es la ausencia de términos ARMA estacionales,
excluidos por su coste de cálculo (sección~\ref{sec:esp-sarimax}), y se
declara como limitación. Con el diagnóstico de esta sección se cumple la
parte del criterio del objetivo~OE3 referida a la especificación final y al
contraste de Ljung-Box; la otra parte, las previsiones del periodo de
prueba, corresponde a la evaluación de 2024.

% ---------------------------------------------------------------------
\subsection{Comparación de los cuatro modelos}\label{sec:val-comparacion}

La tabla~\ref{tab:validacion} reúne las métricas de los cuatro modelos
seleccionados en 2023, calculadas con el programa común de la
sección~\ref{sec:metricas-puntuales}, que reproduce las cifras obtenidas en
la búsqueda de cada familia.

\begin{table}[htbp]
	\centering
	\caption{Métricas de los cuatro modelos en 2023 (validación)}
	\label{tab:validacion}
	\small
	\setlength{\tabcolsep}{3.5pt}
	\begin{tabular}{lrrrrrrrr}
		\hline
		\textbf{Modelo} & \textbf{MAE} & \textbf{RMSE} & \textbf{MAPE} & \textbf{MASE} & \textbf{MAE rel.} & \textbf{Sesgo} & \textbf{RMSE/MAE} & \textbf{P95 $|e|$} \\
		 & (MW) & (MW) & (\%) & & & (MW) & & (MW) \\
		\hline
		Ingenuo & 1400,0 & 2014,9 & 5,32 & 1,054 & 1,000 & $+46$ & 1,44 & 4307 \\
		SARIMAX & 943,4 & 1414,4 & 3,67 & 0,710 & 0,674 & $-346$ & 1,50 & 2767 \\
		LightGBM & 585,2 & 790,7 & 2,25 & 0,441 & 0,418 & $-142$ & 1,35 & 1617 \\
		N-HiTS & 716,4 & 926,5 & 2,73 & 0,539 & 0,512 & $+188$ & 1,29 & 1855 \\
		\hline
	\end{tabular}
\end{table}

Los tres modelos ajustados mejoran al ingenuo estacional, que falla de media
1400~MW por hora, un 5,3\,\% de la demanda, y un 5,4\,\% más que en su
propio periodo de referencia (MASE de 1,054). LightGBM elimina el 58\,\% de
ese error, N-HiTS el 49\,\% y el SARIMAX el 33\,\%. El orden es el mismo en
las cuatro métricas de precisión. El MAE de LightGBM es un 38\,\% menor que el
del SARIMAX y un 18\,\% menor que el de N-HiTS, y el de N-HiTS, un 24\,\%
menor que el del SARIMAX. En 2023, por tanto, el
modelo de aprendizaje profundo no supera al de \textit{gradient boosting}, un
resultado que se somete a la prueba de 2024 y se retoma en la discusión.

Las dos últimas columnas describen cómo se reparte el error. El cociente
entre el RMSE y el MAE, que valdría 1,25 con errores normales centrados, es
máximo en el
SARIMAX (1,50) y en el ingenuo (1,44): una parte importante de su error se
concentra en unas pocas horas con errores muy grandes, que la
sección~\ref{sec:val-dias} localiza. N-HiTS tiene el cociente más bajo, 1,29,
y la menor distancia relativa a LightGBM en el percentil 95 del error absoluto:
su MAE es un 22\,\% mayor que el de LightGBM, pero su percentil 95 solo un
15\,\% mayor. Los sesgos tienen signos distintos: el SARIMAX y LightGBM
prevén de media más demanda de la que hubo, y N-HiTS menos.

% ---------------------------------------------------------------------
\subsection{El error a lo largo del día y del año}\label{sec:val-perfil}

La figura~\ref{fig:val-anticipacion} muestra el MAE de cada modelo según la
anticipación de la hora prevista, medida desde el origen, con la hora local
del día $D+1$ correspondiente en el eje superior. Con un único origen diario,
la anticipación y la hora del día van unidas: la anticipación de 12 horas es
siempre la medianoche, y la de 26 horas, siempre las dos de la tarde, salvo
en los días de cambio de hora. Por eso la figura no permite separar el efecto
de prever más lejos del de prever una hora del día más o menos difícil. Lo que
muestra es que el error no crece de forma regular con la anticipación, sino
que sigue la hora: las horas de madrugada, a 12-17 horas del origen, son las
más fáciles para los cuatro modelos, y las de la tarde, las más difíciles.

\begin{figure}[htbp]
	\centering
	\includegraphics[width=\textwidth]{validacion_mae_anticipacion.png}
	\caption{MAE de los cuatro modelos en 2023 según la anticipación de la hora prevista (eje inferior) y la hora local del día $D+1$ correspondiente (eje superior). Se omite la anticipación de 36~h, que solo se da en una hora del año.}
	\label{fig:val-anticipacion}
\end{figure}

El ingenuo, el SARIMAX y LightGBM tienen el mismo perfil: error bajo de
madrugada, subida a partir de las siete de la mañana y máximo entre las 12 y
las 17 horas, que en LightGBM llega a unos 860~MW. N-HiTS tiene un perfil
propio, en dientes de sierra, con máximos de error a las 07:00 (976~MW) y a
las 21:00 (1315~MW), y es el mejor de los cuatro entre las 13 y las 17 horas.

Para entender ese perfil, la tabla~\ref{tab:perfil} compara la demanda media
real de cada hora de 2023 con la previsión media de cada modelo para esa hora.
El SARIMAX y LightGBM reproducen la forma del día casi exactamente: la
diferencia entre su hora más alta y su hora más baja, la amplitud del perfil,
es del 101-102\,\% de la real, y sus errores medios son desplazamientos de
nivel, no deformaciones. N-HiTS, en cambio, dibuja un día aplanado: su
amplitud es solo el 88\,\% de la real, sube tarde por la mañana, se queda
1280~MW por debajo del pico de las 21:00 y sitúa su máximo una hora antes. Su
sesgo positivo procede de esa forma: es el saldo de quedarse muy corto por la
mañana y por la tarde y pasarse al mediodía, y no un nivel general demasiado
bajo. El patrón es el mismo en los doce meses y en las tres semillas. Su
sesgo medio por franjas es de $+31$~MW de 00:00 a 06:00, $+262$~MW de 06:00 a
12:00, $-313$~MW de 12:00 a 18:00 y $+772$~MW de 18:00 a 24:00.

\begin{table}[htbp]
	\centering
	\caption{Demanda media real y previsión media de cada modelo en horas seleccionadas de 2023 (MW)}
	\label{tab:perfil}
	\small
	\begin{tabular}{lrrrr}
		\hline
		\textbf{Hora local} & \textbf{Real} & \textbf{SARIMAX} & \textbf{LightGBM} & \textbf{N-HiTS} \\
		\hline
		04:00 (valle) & 21\,060 & 21\,262 & 21\,053 & 21\,145 \\
		07:00 (subida) & 25\,337 & 25\,547 & 25\,236 & 24\,487 \\
		12:00 (mediodía) & 27\,854 & 28\,399 & 28\,240 & 28\,588 \\
		16:00 & 27\,185 & 27\,668 & 27\,618 & 27\,347 \\
		21:00 (pico) & 30\,047 & 30\,403 & 30\,139 & 28\,767 \\
		\hline
		Amplitud del perfil & 8987 & 9141 & 9086 & 7909 \\
		Amplitud relativa a la real & 1,000 & 1,017 & 1,011 & 0,880 \\
		\hline
	\end{tabular}
\end{table}

Un perfil de error tan marcado podía deberse a un desfase horario en el
procesamiento de las previsiones, que la prueba de fuga no detectaría. La
comprobación de alineación de la sección~\ref{sec:prot-fuga} lo descarta.
Calculado sobre las 8754 horas comunes a todos los desplazamientos, que
excluyen las tres primeras y las tres últimas del año, el MAE de N-HiTS es de
716,5~MW comparado con la demanda de la misma hora, y de 1088,4 y 1164,3~MW
comparado con la de una hora antes y una hora después; los otros tres
modelos se comportan igual. El aplanamiento es, por tanto, un
comportamiento del modelo. Una explicación compatible con su arquitectura
(sección~\ref{sec:nhits}) es que las pilas de baja resolución dibujan una
trayectoria suave y que la red, que llegó al máximo de pasos de entrenamiento,
no llegó a añadir todo el detalle de la forma diaria. No se ha comprobado,
y el caso es distinto del de las enmiendas del SARIMAX. N-HiTS ya recibe toda
la información que el protocolo pone a disposición de los modelos: la ventana
de demanda hasta el origen y las exógenas. Corregir su perfil exigiría
cambiar hiperparámetros o el procedimiento de entrenamiento, como el número
de pasos, los pesos del horizonte o las escalas, y eso, según la regla de la
sección~\ref{sec:prot-principio}, sería decidir con 2023.

La tabla muestra también un rasgo común a los tres modelos ajustados: los
tres prevén al mediodía más demanda de la que hubo. Entre las 12:00 y las
18:00, su sesgo medio es de $-508$~MW en el SARIMAX, $-411$~MW en LightGBM y
$-313$~MW en N-HiTS, mientras que el del ingenuo, que copia la semana
anterior, es de $+53$~MW. En LightGBM, ese exceso aparece en diez de los doce
meses y es mayor de abril a septiembre, con $-507$~MW de media, que de
noviembre a febrero, con $-210$~MW. La diferencia, de 297~MW, supera el umbral
de 150~MW fijado antes de calcularla para considerarla apreciable. Es
compatible con el efecto del autoconsumo, que reduce la demanda observada en
las horas de sol (sección~\ref{sec:datos-demanda}), pero no lo demuestra, y su
interpretación se aborda en la discusión.

\begin{figure}[htbp]
	\centering
	\includegraphics[width=\textwidth]{validacion_mae_mes.png}
	\caption{MAE mensual de los cuatro modelos en 2023.}
	\label{fig:val-mes}
\end{figure}

La figura~\ref{fig:val-mes} muestra el error de cada mes. LightGBM es el
mejor en diez de los doce meses; N-HiTS lo supera en abril y en agosto. El
ingenuo tiene sus peores meses en agosto (2399~MW), enero y diciembre, los
meses en que la semana anterior se parece menos a la que se prevé, por las
vacaciones y la Navidad. Agosto es también el peor mes del SARIMAX y de
LightGBM. N-HiTS es el más regular a lo largo del año, con errores mensuales
entre 587 y 852~MW, y no empeora en agosto como los demás.

% ---------------------------------------------------------------------
\subsection{Días difíciles y relación entre los modelos}\label{sec:val-dias}

La tabla~\ref{tab:tipo-dia-val} desglosa el error por tipo de día. Los diez
festivos de 2023 concentran los errores más grandes del ingenuo y del
SARIMAX, con un MAE de más de 4000~MW en ambos y un sesgo casi igual al error:
los dos prevén un día laborable cuando la demanda es la de un domingo. Que el
SARIMAX falle en los festivos casi tanto como el ingenuo indica que la
fracción de población en festivo, que entra en su regresión con un efecto
lineal y fijo, no basta para recoger cómo cambia el perfil de un festivo.
Ocho de sus diez peores días son festivos, lo que explica su elevado cociente
entre el RMSE y el MAE. LightGBM y N-HiTS reducen mucho ese error, y N-HiTS es
el mejor de los cuatro en festivos y en puentes. En los cuatro modelos, el
sesgo en festivos es negativo.

\begin{table}[htbp]
	\centering
	\caption{MAE por tipo de día en 2023 (MW)}
	\label{tab:tipo-dia-val}
	\small
	\begin{tabular}{lrrrrr}
		\hline
		\textbf{Tipo de día} & \textbf{Días} & \textbf{Ingenuo} & \textbf{SARIMAX} & \textbf{LightGBM} & \textbf{N-HiTS} \\
		\hline
		Laborable & 247 & 1431 & 882 & 570 & 714 \\
		Fin de semana & 105 & 1045 & 773 & 561 & 707 \\
		Festivo & 10 & 4260 & 4194 & 1131 & 910 \\
		Puente & 3 & 1712 & 1083 & 916 & 567 \\
		\hline
	\end{tabular}
\end{table}

Los peores días de cada modelo reflejan sus debilidades propias. Los del
ingenuo son los festivos y, sobre todo, los días situados una semana después
de un festivo: el 13 de enero, que copia el 6, el 14 de abril, que copia el
Viernes Santo, y del mismo modo el 8 de mayo, el 22 de agosto, el 19 de
octubre y el 8 de noviembre. En todos ellos el modelo prevé un festivo en un
día laborable, y se queda corto en más de 4800~MW de media diaria. Los de LightGBM
son el 1 de mayo, el Viernes Santo, la Navidad y varios días de agosto, en
plenas vacaciones. Entre los de N-HiTS destacan el principio y el final del
año: el 2 y el 3 de enero, el 25 y el 31 de diciembre. Por día de la
semana, N-HiTS se queda especialmente corto los lunes, con un sesgo de
$+490$~MW, cuando su ventana de entrada termina en pleno fin de semana; el
SARIMAX, en cambio, se pasa los lunes ($-524$~MW).

La correlación entre los errores horarios de los modelos indica en qué
medida se equivocan a la vez. La más alta es la del ingenuo con el SARIMAX,
0,70, que fallan en los mismos días; siguen la de LightGBM con el SARIMAX,
0,60, y con el ingenuo, 0,59. N-HiTS es el modelo cuyos errores se parecen
menos a los de los demás: su correlación es de 0,50 con LightGBM, 0,45 con el
SARIMAX y 0,27 con el ingenuo. Por último, el modelo con menor MAE diario es
LightGBM en el 48,8\,\% de los días, N-HiTS en el 22,2\,\%, el SARIMAX en el
16,7\,\% y el ingenuo en el 12,3\,\%: en una semana sin acontecimientos
especiales, repetir la anterior es difícil de mejorar.

% ---------------------------------------------------------------------
\subsection{Síntesis}\label{sec:val-sintesis}

En la validación de 2023, los tres modelos ajustados mejoran con claridad al
ingenuo estacional, y el orden es el mismo en todas las métricas de
precisión: LightGBM, N-HiTS y el SARIMAX. Cada modelo tiene un patrón de error
característico: el SARIMAX, una dependencia diaria residual y un fallo
pronunciado en los festivos; LightGBM, un exceso de previsión al mediodía; y
N-HiTS, un perfil diario aplanado que no alcanza el pico de la tarde, junto
con una variabilidad apreciable entre semillas, de 66~MW, y los errores menos
correlacionados
con los de los demás. La selección se hizo solo con el MAE; los desgloses,
calculados después con reglas de lectura fijadas de antemano, no cambiaron
ninguna ganadora. Si estos resultados se mantienen, y si las diferencias son
significativas, se comprueba en la sección~\ref{sec:evaluacion-final} con la
evaluación única de 2024, donde el orden de los tres modelos ajustados se
mantiene y el SARIMAX deja de distinguirse del ingenuo.

% =====================================================================
% Sección 5.3 redactada el 04/11/2026 (sesión R2). Cifras de
% data/resultados/probabilistica/ generadas por src/27_capa_probabilistica.py
% (protocolo, apartado 14.13). Figuras copiadas de ese directorio.
% =====================================================================
\section{Calibración en validación (2023)}\label{sec:calibracion}

Esta sección presenta la calibración de los intervalos de los cuatro modelos en
2023, con la capa probabilística de la sección~\ref{sec:probabilistica}. Como
en la sección~\ref{sec:validacion}, conviene recordar el papel de estas cifras:
describen cómo se comporta cada método y sirven para elegir el paso $\gamma$ de
ACI, pero no son el rendimiento que se informa. Los hiperparámetros de LightGBM
y de N-HiTS se eligieron con el MAE de 2023, y la ganadora de N-HiTS lleva
incluidos sus cuantiles 0,05 y 0,95, de modo que la calibración de 2023 es
algo optimista; el $\gamma$ de ACI se elige con los mismos datos que se
informan. El rendimiento que cuenta es el de 2024.

La sección comienza con los intervalos propios del 90\,\% de cada modelo,
sigue con la fiabilidad de la rejilla de cuantiles, estudia después ACI y la
elección de $\gamma$, y termina con una síntesis. Los intervalos propios se
evalúan en las 8760 horas del año; la comparación con ACI, que no emite en los
57 primeros días, se hace en el conjunto de 308 días de la
sección~\ref{sec:prob-aci}, al que se llama conjunto~E.

% ---------------------------------------------------------------------
\subsection{Intervalos propios del 90\,\%}\label{sec:cal-propios}

La tabla~\ref{tab:cal-propios} recoge las métricas de los intervalos propios en
las 8760 horas de 2023. Solo el ingenuo está bien calibrado; los otros tres
modelos lo están mal, cada uno a su manera.

\begin{table}[htbp]
	\centering
	\caption{Intervalos propios del 90\,\% en 2023 (8760 horas)}
	\label{tab:cal-propios}
	\small
	\begin{tabular}{lrrrr}
		\hline
		 & \textbf{Ingenuo} & \textbf{SARIMAX} & \textbf{LightGBM} & \textbf{N-HiTS} \\
		\hline
		Cobertura (\%) & 88,3 & 92,5 & 70,0 & 68,1 \\
		\quad IC del 95\,\% (\%) & 85,6--90,8 & 90,4--94,2 & 67,3--72,5 & 66,3--70,0 \\
		Fallos por debajo / por encima (\%) & 5,7 / 5,9 & 4,9 / 2,7 & 18,3 / 11,7 & 13,2 / 18,7 \\
		Cociente de fallos & 0,97 & 1,84 & 1,57 & 0,70 \\
		Anchura media (MW) & 5925 & 4385 & 1926 & 1781 \\
		Puntuación de intervalo (MW) & 9357 & 6679 & 4788 & 4899 \\
		CRPS$_{19}$ (MW) & 1038 & 707 & 446 & 531 \\
		CRPS$_{19}$ / MAE & 0,741 & 0,750 & 0,762 & 0,742 \\
		Desviación máxima de fiabilidad & 0,050 & 0,165 & 0,163 & 0,184 \\
		\hline
	\end{tabular}
\end{table}

El \textbf{ingenuo} es el único bien calibrado, con una cobertura del 88,3\,\%
y fallos repartidos a partes iguales por debajo y por encima. Es lo esperable,
porque sus cuantiles proceden de casi 49\,000 errores reales de seis años; en el
conjunto~E su cobertura es del 88,0\,\%, en el borde de la banda del 88 al
92\,\%. Pero es el listón más ancho, con 5925~MW de media. El \textbf{SARIMAX}
tiene sobrecobertura (92,5\,\% en las 8760 horas y 93,0\,\% en el
conjunto~E, con el intervalo de confianza del 95\,\% por encima del 90\,\%) y
sus fallos son asimétricos: 4,9\,\% por debajo del límite inferior frente al
2,7\,\% por encima, coherente con el sesgo negativo de su previsión
(sección~\ref{sec:val-comparacion}). \textbf{LightGBM} y \textbf{N-HiTS} tienen
una infracobertura muy marcada, del 70,0\,\% y del 68,1\,\% para un 90\,\%
nominal, con intervalos de confianza que quedan a más de 17 puntos del
nominal. Es el resultado central de esta sección, y es negativo: los cuantiles
entrenados con la pérdida de pinball, tal como se obtienen de estos dos modelos,
no están calibrados. Los fallos de LightGBM son algo más frecuentes por debajo
(18,3\,\% frente a 11,7\,\%, cociente 1,57), y los de N-HiTS, por encima
(cociente 0,70). Ninguno de los dos cocientes pasa de 2 ni baja de 0,5, de modo
que, por la regla de lectura de la sección~\ref{sec:prob-lectura}, no se leen
como intervalos desplazados, sino como intervalos estrechos. El cociente del
SARIMAX es 1,84 en las 8760 horas, bajo el umbral de 2, y 2,41 en el
conjunto~E. La lectura fijada de antemano usa las 8760 horas, así que no se
lee como desplazado, aunque en el conjunto~E superaría el umbral.

Las puntuaciones ordenan los modelos igual que la precisión puntual: LightGBM,
N-HiTS, SARIMAX e ingenuo. Todas las diferencias entre pares son apreciables
según la regla de lectura, salvo una: entre LightGBM y N-HiTS, la puntuación de
intervalo no se distingue (diferencia de $-104$~MW, con un intervalo de
confianza de $-607$ a $+441$~MW), aunque el CRPS$_{19}$ sí es menor en
LightGBM ($-84$~MW). Con los intervalos propios, el CRPS$_{19}$ es casi
proporcional al MAE (entre 0,74 y 0,76 veces su valor), de modo que no añade
información sobre el orden de los modelos. Esto tiene una consecuencia para
leer la tabla: el buen CRPS$_{19}$ de LightGBM y de N-HiTS no significa que sus
intervalos estén bien calibrados. Se debe a que sus intervalos son estrechos y
sus errores pequeños; una puntuación que premia la nitidez los favorece aunque
dejen fuera el 30\,\% de las horas.

\paragraph{Cobertura por franja horaria.} La tabla~\ref{tab:cal-bloques} desglosa
la cobertura por bloque en el conjunto~E. Como se explicó en la
sección~\ref{sec:prob-principio}, los bloques son franjas horarias del día
objetivo y, a la vez, tramos de anticipación; las diferencias no se pueden
atribuir a la anticipación.

\begin{table}[htbp]
	\centering
	\caption{Cobertura propia del 90\,\% por bloque en el conjunto~E (\%)}
	\label{tab:cal-bloques}
	\small
	\begin{tabular}{lrrrrr}
		\hline
		 & \textbf{B1} & \textbf{B2} & \textbf{B3} & \textbf{B4} & \textbf{Global} \\
		 & (00--05 h) & (06--11 h) & (12--17 h) & (18--23 h) & \\
		\hline
		Ingenuo & 87,8 & 89,0 & 87,9 & 87,4 & 88,0 \\
		SARIMAX & 99,0 & 92,3 & 88,9 & 91,7 & 93,0 \\
		LightGBM & 67,0 & 73,6 & 67,7 & 75,3 & 70,9 \\
		N-HiTS & 80,2 & 68,7 & 64,1 & 56,1 & 67,3 \\
		\hline
	\end{tabular}
\end{table}

Por la regla de lectura, un bloque se distingue del conjunto solo si difiere en
más de cinco puntos de la cobertura global. Ocurre en dos modelos. En el SARIMAX,
la cobertura de la madrugada (B1) es del 99,0\,\%, seis puntos por encima de la
global. Es una consecuencia de que su anchura es casi plana a lo largo del
día, entre 3898~MW a las 00:00 y 4746~MW a las 23:00, mientras que su error
medio va de 481 a 1277~MW: la banda es demasiado ancha cuando el error es
pequeño, de madrugada, y apenas suficiente cuando es grande. En N-HiTS la
cobertura baja del 80,2\,\% en B1 al 56,1\,\% en B4, y el motivo es su perfil
de sesgos por hora, de $+850$~MW a las 07:00, $-734$ a las 12:00 y $+1279$ a
las 21:00, con bandas de entre 1,5 y 3 veces su error medio; a las 21:00 su
cobertura es del 32\,\% y a las 05:00, del 86\,\%. En LightGBM la banda sigue la
forma del día, de 2,6 a 5 veces su error medio, pero es demasiado estrecha a
todas horas, y ningún bloque se aparta más de cinco puntos de la cobertura
global. En el ingenuo no se distingue ninguno.

La figura~\ref{fig:cal-semana} muestra los intervalos en una semana de mayo de
2023. Es un ejemplo y no una prueba: sirve para ver lo que las tablas resumen. La
banda del SARIMAX tiene una anchura casi constante, la de LightGBM sigue la
forma del día pero es estrecha, y la de N-HiTS es la más ajustada.

\begin{figure}[htbp]
	\centering
	\includegraphics[width=\textwidth]{calibracion_semana.png}
	\caption{Previsión, demanda real, intervalo propio del 90\,\% e intervalo de ACI (con $\gamma = 0{,}001$) de los cuatro modelos en la semana del 15 al 21 de mayo de 2023 (días objetivo).}
	\label{fig:cal-semana}
\end{figure}

Por último, la cobertura varía mucho de un mes a otro. Agosto es el mes de
menor cobertura para el ingenuo (67,7\,\%) y para LightGBM (55,9\,\%), y
también baja en el SARIMAX (83,9\,\%). Los datos muestran que el mes es
difícil para todos, pero no se ha comprobado a qué se debe, y su explicación
queda como hipótesis para la discusión.

% ---------------------------------------------------------------------
\subsection{Fiabilidad de la rejilla}\label{sec:cal-fiabilidad}

El diagrama de fiabilidad de la figura~\ref{fig:cal-fiabilidad} generaliza la
cobertura del intervalo del 90\,\% a toda la rejilla de 19 cuantiles. Las curvas
azules son las de los cuantiles propios, evaluados en el conjunto~E; las
naranjas, las de los cuantiles de ACI, que se discuten más adelante.

\begin{figure}[htbp]
	\centering
	\includegraphics[width=\textwidth]{calibracion_fiabilidad.png}
	\caption{Diagrama de fiabilidad de los cuantiles propios y de los de ACI (con $\gamma = 0{,}001$) en el conjunto~E. En el eje horizontal, el nivel del cuantil; en el vertical, la frecuencia con que la demanda queda por debajo de él.}
	\label{fig:cal-fiabilidad}
\end{figure}

La curva del ingenuo está sobre la diagonal. La del SARIMAX queda por encima en
los niveles medios y altos, es decir, sus cuantiles son demasiado altos, que es
la huella de su sesgo negativo: la demanda cae por debajo del cuantil 0,5 el
65\,\% de las veces en el conjunto~E. LightGBM y N-HiTS dibujan una curva más
plana que la diagonal: sus cuantiles bajos son demasiado altos y los altos
demasiado bajos, que es lo que significa un intervalo estrecho. En las 8760
horas, con la rejilla de LightGBM, la demanda cae por debajo del cuantil 0,05
el 16,5\,\% de las veces y por encima del 0,95 el 10,6\,\%, cuando deberían
ser el 5\,\% en ambos casos. Por la regla de lectura, la desviación máxima de los intervalos
propios es \emph{mala} en el SARIMAX, LightGBM y N-HiTS (más de 0,06) y
\emph{aceptable} en el ingenuo (0,050).

Dos detalles técnicos de la rejilla. Los 19 cuantiles de LightGBM y de N-HiTS se
entrenan por separado y pueden cruzarse, de modo que hay que ordenarlos: la
ordenación fue necesaria en las 8760 horas de LightGBM y en 7737 de N-HiTS, y no
lo fue en el ingenuo ni en el SARIMAX. En LightGBM el efecto es pequeño:
los cuantiles 0,05 y 0,95 cambian de media unos 30~MW, frente a una anchura de
unos 1900~MW. La segunda red de N-HiTS, la de los 19 niveles, llega al tope de
1000 pasos; el MAE de su mediana es de 710,3~MW, algo menor que el de la red
oficial (716,4~MW), y la cobertura de sus cuantiles 0,05 y 0,95 es del
66,3\,\%, a menos de dos puntos de la de la red oficial (68,1\,\%). Son, pues,
dos sistemas parecidos, aunque no el mismo.

% ---------------------------------------------------------------------
\subsection{ACI y elección de $\gamma$}\label{sec:cal-aci}

La tabla~\ref{tab:cal-gamma} recoge la elección del paso $\gamma$ de ACI, con la
regla de la sección~\ref{sec:prob-aci}. Cada celda es el cociente entre la
puntuación de intervalo con ese $\gamma$ y la de $\gamma = 0$, la ventana móvil
sin adaptación, en el conjunto~E; la última columna es la media de los cuatro
modelos.

\begin{table}[htbp]
	\centering
	\caption{Puntuación de intervalo de ACI relativa a la de la ventana sin adaptación ($\gamma = 0$), en el conjunto~E}
	\label{tab:cal-gamma}
	\small
	\begin{tabular}{lrrrrr}
		\hline
		$\gamma$ & \textbf{Ingenuo} & \textbf{SARIMAX} & \textbf{LightGBM} & \textbf{N-HiTS} & \textbf{Media} \\
		\hline
		0,001 & 0,998 & 1,019 & 1,001 & 1,001 & 1,0049 \\
		0,005 & 1,026 & 1,059 & 1,010 & 1,006 & 1,0251 \\
		0,01 & 0,996 & 1,052 & 0,990 & 1,013 & 1,0128 \\
		0,02 & 0,989 & 1,038 & 0,989 & 1,021 & 1,0093 \\
		\hline
	\end{tabular}
\end{table}

El resultado es negativo: \textbf{ningún candidato queda por debajo de uno}. La
adaptación de $\alpha_t$ no mejora la puntuación de intervalo de la ventana móvil
sin adaptación. Por la regla, se mantiene $\gamma = 0{,}001$, el de menor
puntuación (1,0049), que es, en la práctica, casi inerte: el nivel $\alpha_t$ medio
de los cuatro modelos queda entre 0,079 y 0,098, cerca del 0,10 inicial. La
razón de que los pasos mayores empeoren se ve en el recuento de intervalos
infinitos sustituidos (tabla~\ref{tab:cal-infinitos}): con $\gamma$ grande, un
día malo puede llevar $\alpha_t$ a cero, porque cada fallo lo reduce en
$0{,}9\gamma$ y la retroalimentación llega en bloques de seis horas, y el
intervalo emitido es entonces la extensión máxima de la ventana. Con
$\gamma = 0{,}02$, el 4,1\,\% de las horas del ingenuo y el 3,9\,\% de las del
SARIMAX tienen límite infinito, por encima del 2\,\% que la regla de lectura
fija como defecto del método. Esta interpretación es coherente con los
recuentos, pero no se ha comprobado por separado.

\begin{table}[htbp]
	\centering
	\caption{Horas con límite infinito sustituido en el conjunto~E (de 7392)}
	\label{tab:cal-infinitos}
	\small
	\begin{tabular}{lrrrr}
		\hline
		$\gamma$ & \textbf{Ingenuo} & \textbf{SARIMAX} & \textbf{LightGBM} & \textbf{N-HiTS} \\
		\hline
		0,001 & 0 & 0 & 0 & 0 \\
		0,005 & 36 & 6 & 6 & 0 \\
		0,01 & 162 & 114 & 96 & 18 \\
		0,02 & 300 & 288 & 216 & 114 \\
		\hline
	\end{tabular}
\end{table}

La tabla~\ref{tab:cal-aci} compara, con $\gamma = 0{,}001$ y en el
conjunto~E, el intervalo propio, la ventana sin adaptación y ACI. Lo
importante es que, para LightGBM y N-HiTS, la mejora procede de la ventana de
residuos recientes y no de la adaptación.

\begin{table}[htbp]
	\centering
	\caption{Intervalo propio, ventana sin adaptación y ACI ($\gamma = 0{,}001$) en el conjunto~E}
	\label{tab:cal-aci}
	\small
	\setlength{\tabcolsep}{3.5pt}
	\begin{tabular}{lrrrr}
		\hline
		 & \textbf{Ingenuo} & \textbf{SARIMAX} & \textbf{LightGBM} & \textbf{N-HiTS} \\
		\hline
		\textit{Cobertura (\%)} & & & & \\
		\quad Propio & 88,0 & 93,0 & 70,9 & 67,3 \\
		\quad Ventana ($\gamma = 0$) & 86,3 & 87,9 & 88,0 & 89,2 \\
		\quad ACI & 88,4 & 89,0 & 89,2 & 89,1 \\
		\quad IC del 95\,\% de ACI & 85,5--91,3 & 86,3--91,6 & 87,0--91,3 & 87,5--90,7 \\
		\textit{Anchura media (MW)} & & & & \\
		\quad Propio & 5925 & 4385 & 1913 & 1727 \\
		\quad Ventana ($\gamma = 0$) & 5924 & 3810 & 2339 & 2506 \\
		\quad ACI & 6568 & 4203 & 2439 & 2514 \\
		\textit{Puntuación de intervalo (MW)} & & & & \\
		\quad Propio & 9446 & 6646 & 4764 & 4868 \\
		\quad Ventana ($\gamma = 0$) & 10\,018 & 6655 & 3324 & 3321 \\
		\quad ACI & 10\,001 & 6785 & 3327 & 3325 \\
		\textit{ACI $-$ propio (MW)} & & & & \\
		\quad Puntuación de intervalo & $+555$ & $+139$ & $-1437$ & $-1544$ \\
		\quad Lectura & {\footnotesize apreciable (peor)} & {\footnotesize no distinguible} & apreciable & apreciable \\
		\quad CRPS$_{19}$ & $+23$ & $-40$ & $-35$ & $-87$ \\
		\quad Lectura & {\footnotesize no distinguible} & apreciable & apreciable & apreciable \\
		\textit{Fiabilidad de ACI} & & & & \\
		\quad Desviación máxima & 0,039 & 0,009 & 0,031 & 0,010 \\
		\hline
	\end{tabular}
\end{table}

En \textbf{LightGBM} y \textbf{N-HiTS}, ACI lleva la cobertura del 70,9\,\% y del
67,3\,\% a cerca del 89\,\% y baja la puntuación de intervalo un 30\,\% y un
32\,\% respecto a la del intervalo propio, con diferencias apreciables. Pero la
ventana sin adaptación ya consigue casi lo mismo (puntuaciones de 3324 y
3321~MW frente a 3327 y 3325~MW con ACI): la mejora viene de calcular el
intervalo con los residuos de las últimas ocho semanas, no de la adaptación de
$\alpha_t$. En el \textbf{SARIMAX}, ACI acerca la cobertura al nominal (la
distancia al 90\,\% pasa de 0,030 a 0,010) y mejora el CRPS$_{19}$, pero no la
puntuación de intervalo, que no se distingue. En el \textbf{ingenuo}, ACI
empeora de forma apreciable la puntuación de intervalo ($+555$~MW) y no mejora
la calibración (la distancia pasa de 0,020 a 0,016, menos de un punto): una
ventana de 336 residuos pierde frente a los 49\,320 errores de seis años con que
se construye su intervalo propio. La cobertura con ACI sigue oscilando mucho de
un mes a otro; la del ingenuo, por ejemplo, va del 81\,\% al 100\,\%.

Conviene repetir lo que ACI es en este trabajo (sección~\ref{sec:prob-aci}):
un intervalo construido con los residuos recientes, que no usa el intervalo
propio. Por eso su comparación con él no mide el efecto de una recalibración del
intervalo de cada modelo, sino el de cambiar de mecanismo. Y, aun así, el
resultado de este cambio es muy distinto según el modelo de partida: no aporta
nada a quien ya estaba calibrado, como el ingenuo, y arregla en gran parte a
quienes no lo estaban.

% ---------------------------------------------------------------------
\subsection{Síntesis}\label{sec:cal-sintesis}

La calibración de 2023 deja cuatro resultados. Primero, los cuantiles de
LightGBM y de N-HiTS entrenados con la pérdida de pinball no están calibrados,
con coberturas del 70\,\% y del 68\,\% para un 90\,\% nominal, y su buen
CRPS$_{19}$ procede de intervalos estrechos, no de intervalos fiables. Segundo,
el intervalo gaussiano homocedástico del SARIMAX cubre de más y reparte mal su
anchura a lo largo del día: demasiado ancho de madrugada y justo por la tarde.
Tercero, calcular el intervalo con una ventana de residuos recientes, con o sin
ACI, lleva a ambos modelos a cerca del 90\,\% y reduce un 30\,\% su
puntuación de intervalo, pero la adaptación de ACI no aporta nada sobre la
ventana, y el $\gamma$ elegido es casi inerte. Cuarto, el ingenuo, con su
intervalo propio basado en errores de seis años, es el único método bien
calibrado sin ayuda, aunque a costa de los intervalos más anchos.

Estos resultados son optimistas en el sentido ya señalado: los hiperparámetros y
el $\gamma$ se eligieron con 2023. Para 2024 se mantienen el
$\gamma = 0{,}001$ y las mismas reglas de lectura, con la ventana y ACI
evaluados sobre las 8784 horas del año.

% =====================================================================
\section{Evaluación final en 2024: precisión puntual}\label{sec:evaluacion-final}

Esta sección presenta el rendimiento que se informa del trabajo. 2024 es el
año que no intervino en ninguna elección: las cuatro configuraciones eran las
ganadoras de 2023, se congelaron antes de calcular nada, los modelos se
reajustaron cada mes con todo lo ocurrido hasta el origen
(sección~\ref{sec:prot-particiones}) y las cifras se leyeron una sola vez, con las
reglas de la sección~\ref{sec:contraste-dm}. Las previsiones tienen una
anticipación de 12 a 36 horas y cubren las 8784 horas del año bisiesto.

La sección empieza con la precisión de los cuatro modelos, sigue con el
contraste de Diebold--Mariano, compara después con 2023, describe el error por
mes y por franja horaria y termina con una síntesis.

\subsection{Precisión de los cuatro modelos}\label{sec:final-precision}

La tabla~\ref{tab:final-puntuales} recoge las métricas puntuales. El orden por
MAE es el mismo que en la validación: LightGBM (573,6~MW), N-HiTS (665,6),
SARIMAX (1089,7) e ingenuo (1217,9). Medido contra el ingenuo, LightGBM
elimina el 53\,\% de su error absoluto y N-HiTS el 45\,\%, mientras que el
SARIMAX solo elimina el 10,5\,\%. El MASE del ingenuo vale 0,917 y no uno
porque la escala es fija (sección~\ref{sec:metricas-puntuales}): 2024 fue un año
algo más fácil para el ingenuo que el periodo de 2017 a 2022 con que se fijó
la escala.

\begin{table}[htbp]
	\centering
	\caption{Precisión puntual de los cuatro modelos en 2024 (8784 horas)}
	\label{tab:final-puntuales}
	\small
	\begin{tabular}{lrrrr}
		\hline
		 & \textbf{Ingenuo} & \textbf{SARIMAX} & \textbf{LightGBM} & \textbf{N-HiTS} \\
		\hline
		MAE (MW) & 1217,9 & 1089,7 & 573,6 & 665,6 \\
		\quad IC del 95\,\% (MW) & 1107--1342 & 984--1195 & 545--605 & 633--700 \\
		MAE relativo al ingenuo & 1,000 & 0,895 & 0,471 & 0,547 \\
		RMSE (MW) & 1871,7 & 1616,0 & 770,7 & 875,8 \\
		MAPE (\%) & 4,56 & 4,18 & 2,16 & 2,54 \\
		MASE & 0,917 & 0,820 & 0,432 & 0,501 \\
		Sesgo (MW) & $-$6,7 & $+$6,0 & $+$135,0 & $-$41,1 \\
		RMSE / MAE & 1,54 & 1,48 & 1,34 & 1,32 \\
		\hline
	\end{tabular}
\end{table}


Los intervalos de confianza del 95\,\% del MAE, calculados por
\textit{bootstrap} de días, no se solapan entre LightGBM (545--605~MW) y
N-HiTS (633--700), ni entre estos dos y los clásicos. Sí se solapan los del
SARIMAX (984--1195) y el ingenuo (1107--1342). El solapamiento de dos
intervalos no equivale a un contraste, que se hace en el apartado siguiente,
pero anticipa su resultado. Ningún modelo tiene un sesgo apreciable, es decir,
superior a 150~MW en valor absoluto; el mayor es el de LightGBM, de $+135$~MW,
lo que significa que se queda corto de media (el error es real menos previsto).
La razón entre el RMSE y el MAE es menor en los dos modelos de aprendizaje
automático (1,34 y 1,32) que en los clásicos (1,48 y 1,54), de modo que
sus errores grandes pesan menos.

\subsection{Contraste de Diebold--Mariano}\label{sec:final-dm}

La tabla~\ref{tab:final-dm} contrasta los seis pares con la pérdida absoluta
diaria, el método de Newey y West con 7 retardos, la corrección de Harvey,
Leybourne y Newbold y el ajuste de Holm. Una diferencia media positiva en la
tabla significa que el segundo modelo del par se equivoca menos.

\begin{table}[htbp]
	\centering
	\caption{Contraste de Diebold--Mariano en 2024 (pérdida absoluta diaria, 366 días, 7 retardos)}
	\label{tab:final-dm}
	\small
	\begin{tabular}{llrrrl}
		\hline
		\textbf{Modelo $a$} & \textbf{Modelo $b$} & \textbf{$\bar d$ (MW)} & \textbf{DM} & \textbf{p ajustado} & \textbf{Lectura} \\
		\hline
		ingenuo & SARIMAX & $+$128 & 1,72 & 0,087 & no distinguible \\
		ingenuo & LightGBM & $+$644 & 6,42 & $<0{,}001$ & apreciable \\
		ingenuo & N-HiTS & $+$552 & 5,34 & $<0{,}001$ & apreciable \\
		SARIMAX & LightGBM & $+$516 & 6,81 & $<0{,}001$ & apreciable \\
		SARIMAX & N-HiTS & $+$424 & 5,67 & $<0{,}001$ & apreciable \\
		LightGBM & N-HiTS & $-$92 & $-$4,55 & $<0{,}001$ & apreciable \\
		\hline
	\end{tabular}
\end{table}


Cinco de los seis pares son apreciables, todos con un p-valor ajustado
inferior a 0,001: LightGBM y N-HiTS superan al ingenuo y al SARIMAX, y LightGBM
supera a N-HiTS en 92~MW de pérdida media diaria. La conclusión se mantiene con
la pérdida cuadrática y con 3 y 14 retardos.

El sexto par es un \textbf{resultado negativo}: el SARIMAX no se distingue del
ingenuo estacional. La diferencia a su favor es de 128~MW por día, con un
p-valor ajustado de 0,087. La conclusión depende algo de las decisiones del
contraste. Con 3, 7 y 14 retardos, el p-valor ajustado es 0,057, 0,087 y
0,120. Con pérdida cuadrática es 0,063 con 7 retardos y 0,086 con 14, pero
baja a 0,047 con 3, justo por debajo del umbral. La lectura que se fijó antes
de calcular, con pérdida absoluta y 7 retardos, es \enquote{no distinguible} y
es la que se informa. Lo que dicen las variantes es que, en el mejor caso, el
SARIMAX mejora al ingenuo por un margen muy pequeño. Con la regla de la
sección~\ref{sec:contraste-dm}, el orden de los tres primeros modelos
(LightGBM, N-HiTS y SARIMAX) queda establecido, pero no el de los cuatro,
porque el último par consecutivo no es apreciable.

\subsection{Comparación con 2023}\label{sec:final-2023}

La comparación con la validación es solo descriptiva. El MAE del ingenuo baja un
13,0\,\% (de 1400,0 a 1217,9~MW), el del SARIMAX sube un 15,5\,\% (de 943,4 a
1089,7), el de LightGBM baja un 2,0\,\% (de 585,2 a 573,6) y el de N-HiTS, un
7,1\,\% (de 716,4 a 665,6). Con la regla de lectura, el ingenuo es mejor que en
2023, el SARIMAX es peor y los dos modelos de aprendizaje automático son
similares.

Estos cambios admiten varias explicaciones que el diseño no permite separar.
Una es el efecto del año: el ingenuo mejora un 13\,\% sin que haya cambiado
nada en él, de modo que 2024 es un año distinto para quien copia el patrón de la
semana anterior. Otra es la \enquote{maldición del ganador}, porque los
hiperparámetros de los tres modelos ajustados se eligieron con 2023, y la tercera, el
reajuste mensual. En LightGBM y N-HiTS no se observa una pérdida de más del 10\,\%, pero no
se puede descartar una pequeña. El
SARIMAX es el caso más llamativo: con el reajuste mensual su sesgo pasa de
$-346$~MW en 2023 a $+6$~MW en 2024 y deja de ser apreciable, pero el error no
mejora, sino que empeora. No se ha comprobado a qué se debe. Tampoco se ha
medido el efecto del reajuste por separado, porque no se hizo un control sin
reajuste. Los sesgos de LightGBM y de N-HiTS cambian además de signo (de
$-142$ a $+135$~MW y de $+188$ a $-41$~MW), sin una explicación comprobada.

\subsection{Error por mes y por franja horaria}\label{sec:final-mes}

La tabla~\ref{tab:final-mes} desglosa el MAE por mes y por franja horaria del
día objetivo. Es una descripción, sin contrastes y sin explicación comprobada;
el desglose por hora del día y tipo de día, que sí contrasta las diferencias,
es el de la sección~\ref{sec:desglose}.

\begin{table}[htbp]
	\centering
	\caption{MAE por mes y por franja horaria en 2024 (MW)}
	\label{tab:final-mes}
	\small
	\begin{tabular}{lrrrr}
		\hline
		 & \textbf{Ingenuo} & \textbf{SARIMAX} & \textbf{LightGBM} & \textbf{N-HiTS} \\
		\hline
		Enero & 1936 & 1655 & 623 & 797 \\
		Febrero & 898 & 689 & 669 & 702 \\
		Marzo & 1273 & 1282 & 680 & 866 \\
		Abril & 1042 & 1441 & 745 & 897 \\
		Mayo & 906 & 1048 & 426 & 485 \\
		Junio & 935 & 767 & 562 & 582 \\
		Julio & 1017 & 780 & 518 & 490 \\
		Agosto & 1698 & 988 & 548 & 499 \\
		Septiembre & 985 & 804 & 536 & 662 \\
		Octubre & 723 & 691 & 571 & 742 \\
		Noviembre & 1010 & 1385 & 445 & 556 \\
		Diciembre & 2143 & 1523 & 566 & 713 \\
		\hline
		B1 (00--05 h) & 745 & 825 & 375 & 537 \\
		B2 (06--11 h) & 1265 & 1118 & 578 & 678 \\
		B3 (12--17 h) & 1649 & 1304 & 792 & 762 \\
		B4 (18--23 h) & 1212 & 1112 & 549 & 685 \\
		\hline
	\end{tabular}
\end{table}


Los meses más difíciles son diciembre y enero para el ingenuo (2143 y
1936~MW), enero, diciembre, abril y noviembre para el SARIMAX (de 1385 a
1655~MW) y abril para LightGBM y N-HiTS (745 y 897~MW); el mejor mes de estos
dos es mayo (426 y 485~MW). El SARIMAX queda
por encima del ingenuo en marzo, abril, mayo y noviembre, y en la franja de
madrugada (B1, 825 frente a 745~MW). N-HiTS queda por debajo de LightGBM en
julio (490 frente a 518~MW), en agosto (499 frente a 548) y en la franja de la
tarde (B3, 762 frente a 792). Las diferencias entre los dos modelos de
aprendizaje automático son pequeñas en esos casos y no se han contrastado
aquí. Los cuatro modelos tienen su peor franja entre las 12 y las 17 horas (B3).

\subsection{Síntesis}\label{sec:final-sintesis}

La evaluación final deja cuatro resultados. Primero, el orden de precisión de
2024 es el de la validación: LightGBM, N-HiTS, SARIMAX e ingenuo, y LightGBM es
apreciablemente mejor que N-HiTS, con una diferencia robusta a la pérdida y al
número de retardos. Segundo, y es el resultado negativo, el SARIMAX con
estacionalidad múltiple no se distingue del ingenuo estacional con el criterio
fijado de antemano: el modelo clásico más elaborado del trabajo no mejora de
forma demostrable a la referencia más simple. Tercero, los dos modelos de
aprendizaje automático conservan en 2024 lo que mostraron en 2023, y el
SARIMAX no. Cuarto, el sesgo de ningún modelo es apreciable, pero la comparación
con 2023 mezcla el efecto del año, el del reajuste y el de haber elegido los
hiperparámetros con 2023, y no permite repartir las causas.

% =====================================================================
\section{Calibración en 2024}\label{sec:calibracion-final}

Esta sección presenta la calibración de los intervalos de los cuatro modelos en
2024, el año de la evaluación final. A diferencia de la
sección~\ref{sec:calibracion}, todos los intervalos se evalúan sobre las mismas
8784 horas, porque ACI parte de los residuos de 2023 y emite desde el primer día
del año (sección~\ref{sec:prob-aci}). Las reglas de lectura son las de la
sección~\ref{sec:prob-lectura}, y el paso $\gamma = 0{,}001$ de ACI es el
elegido con 2023: no se reelige con 2024.

La sección sigue el orden de la~\ref{sec:calibracion}: intervalos propios,
fiabilidad de la rejilla, ventana y ACI, y una síntesis.

\subsection{Intervalos propios del 90\,\%}\label{sec:final-propios}

La tabla~\ref{tab:final-propios} recoge las métricas de los intervalos propios.
Los dos modelos clásicos sobrecubren por poco y los dos de aprendizaje automático
infracubren, como en 2023, aunque menos.

\begin{table}[htbp]
	\centering
	\caption{Intervalos propios del 90\,\% en 2024 (8784 horas)}
	\label{tab:final-propios}
	\small
	\begin{tabular}{lrrrr}
		\hline
		 & \textbf{Ingenuo} & \textbf{SARIMAX} & \textbf{LightGBM} & \textbf{N-HiTS} \\
		\hline
		Cobertura (\%) & 92,3 & 93,3 & 76,1 & 78,7 \\
		\quad IC del 95\,\% (\%) & 90,0--94,4 & 91,3--95,2 & 73,7--78,4 & 76,7--80,8 \\
		Fallos por debajo / por encima (\%) & 3,8 / 3,9 & 3,5 / 3,2 & 10,6 / 13,3 & 11,6 / 9,6 \\
		Cociente de fallos & 0,97 & 1,09 & 0,79 & 1,21 \\
		Anchura media (MW) & 5948 & 5717 & 2221 & 2081 \\
		Puntuación de intervalo (MW) & 8811 & 7392 & 4418 & 4157 \\
		CRPS$_{19}$ (MW) & 913 & 826 & 420 & 485 \\
		CRPS$_{19}$ / MAE & 0,749 & 0,758 & 0,732 & 0,729 \\
		Desviación máxima de fiabilidad & 0,023 & 0,096 & 0,141 & 0,084 \\
		\hline
	\end{tabular}
\end{table}


El \textbf{ingenuo} cubre el 92,3\,\%, algo por encima de la banda del 88 al
92\,\% de la regla de lectura; se lee como sobrecobertura, pero por
0,3 puntos y con un intervalo de confianza (90,0--94,4\,\%) que contiene el
90\,\%, de modo que es un resultado marginal. Sus fallos están repartidos a
partes iguales. En 2023 cubría el 88,3\,\%, con cuantiles de 2017--2022, y en
2024, con los cuantiles actualizados cada mes, el 92,3\,\%. El \textbf{SARIMAX}
cubre el 93,3\,\% (91,3--95,2\,\%), con el intervalo de confianza completo por
encima del 90\,\%: sobrecobertura, con una asimetría pequeña (cociente de
fallos 1,09, dentro de la banda de 0,5 a 2). En 2023 era del 92,5\,\%.

\textbf{LightGBM} y \textbf{N-HiTS} cubren el 76,1\,\% (73,7--78,4\,\%) y el
78,7\,\% (76,7--80,8\,\%): la infracobertura de 2023 (70,0 y 68,1\,\%) se
repite, algo menor. Por la regla de lectura, la infracobertura \enquote{se
repite} porque ambas coberturas están por debajo del 88\,\% y sus intervalos de
confianza quedan por debajo del 90\,\%. Es compatible con la idea de que los cuantiles de un modelo de
cuantiles se ajustan mejor en entrenamiento, donde el error es menor, que
fuera de muestra (sección~\ref{sec:conforme}), pero no la demuestra. Ni el cociente de
fallos de LightGBM (0,79) ni el de N-HiTS (1,21) pasan de 2 ni bajan de 0,5, y
por tanto los intervalos se leen como estrechos y no como desplazados.

La puntuación de intervalo ordena los modelos como el MAE, salvo que N-HiTS
queda por delante de LightGBM (4157 frente a 4418~MW). Esa diferencia no es
apreciable: el intervalo de confianza de la diferencia pareada, de $+262$~MW, va
de $-146$ a $+675$~MW. El CRPS$_{19}$, en cambio, sí es apreciablemente menor en
LightGBM (420 frente a 485~MW; diferencia de $-65$~MW, de $-93$ a $-40$). El
resto de las diferencias entre pares son apreciables en las dos puntuaciones. Como en
2023, el CRPS$_{19}$ es casi proporcional al MAE (entre 0,73 y 0,76 veces su valor),
de modo que no añade información sobre el orden de los modelos, y el buen
CRPS$_{19}$ de las dos redes y de LightGBM no significa que sus intervalos estén
calibrados.

\paragraph{Cobertura por franja horaria.} Con la regla de que un bloque se
distingue del conjunto solo si difiere en más de cinco puntos de la cobertura
global, solo hay una franja destacable: la de la tarde (B3) en LightGBM, con el
68,9\,\%, 7,2 puntos por debajo del 76,1\,\% global. Los demás bloques de los
cuatro modelos quedan a menos de cinco puntos de su global. La
figura~\ref{fig:final-bloque-mes} muestra la cobertura por franja y por mes con
el intervalo propio y con el de ACI.

\begin{figure}[htbp]
	\centering
	\includegraphics[width=\textwidth]{calibracion_2024_bloque_mes.png}
	\caption{Cobertura del intervalo del 90\,\% por franja horaria del día objetivo y por mes en 2024, con el intervalo propio (línea continua) y el de ACI con $\gamma = 0{,}001$ (línea discontinua). La banda gris es el 88--92\,\% de la regla de lectura.}
	\label{fig:final-bloque-mes}
\end{figure}

Por mes, la cobertura oscila mucho. La del ingenuo baja al 74,2\,\% en
diciembre, al 79,4\,\% en enero y al 84,8\,\% en agosto, que son sus tres meses
de mayor error absoluto (tabla~\ref{tab:final-mes}), y llega al 100\,\% en julio
y en octubre. Es lo esperable si sus cuantiles, que dependen solo del bloque
horario y de los errores pasados, no distinguen unos meses de otros; pero esta
coincidencia es una descripción, no una prueba. El mínimo mensual de LightGBM es el 62,2\,\% en enero
y el de N-HiTS el 62,2\,\% en abril, su mes de mayor error; el del SARIMAX, el
87,4\,\% en diciembre.

La figura~\ref{fig:final-semana} muestra los intervalos en una semana de mayo
de 2024, la misma posición semanal que la semana de 2023. Es un ejemplo y no una
prueba.

\begin{figure}[htbp]
	\centering
	\includegraphics[width=\textwidth]{calibracion_2024_semana.png}
	\caption{Previsión, demanda real, intervalo propio del 90\,\% e intervalo de ACI (con $\gamma = 0{,}001$) de los cuatro modelos en la semana del 13 al 19 de mayo de 2024 (días objetivo).}
	\label{fig:final-semana}
\end{figure}

\subsection{Fiabilidad de la rejilla}\label{sec:final-fiabilidad}

La figura~\ref{fig:final-fiabilidad} es el diagrama de fiabilidad de la rejilla
de 19 cuantiles, con el intervalo propio y con ACI.

\begin{figure}[htbp]
	\centering
	\includegraphics[width=\textwidth]{calibracion_2024_fiabilidad.png}
	\caption{Diagrama de fiabilidad de la rejilla de cuantiles en 2024: frecuencia con que la demanda queda por debajo del cuantil previsto, para el intervalo propio y con ACI ($\gamma = 0{,}001$).}
	\label{fig:final-fiabilidad}
\end{figure}

Las desviaciones máximas de los modelos propios son 0,023 (ingenuo), 0,096
(SARIMAX), 0,141 (LightGBM) y 0,084 (N-HiTS), que por la regla de lectura son
\emph{buena}, \emph{mala}, \emph{mala} y \emph{mala} (hasta 0,03, buena; hasta
0,06, aceptable; más, mala). El ingenuo sigue la diagonal. El SARIMAX dibuja una
S: la frecuencia queda por debajo de la diagonal en los niveles bajos y por
encima en los altos, lo que corresponde a una distribución prevista más ancha
que la real. En LightGBM la frecuencia observada en el cuantil 0,5 es 0,394:
la demanda real supera a la mediana prevista el 61\,\% de las horas, coherente con
el sesgo de $+135$~MW. En N-HiTS la mediana está centrada (0,524), pero las colas
son demasiado estrechas: la frecuencia por debajo del cuantil 0,05 es 0,132 y
por debajo del 0,90, 0,842.

\subsection{Ventana y ACI}\label{sec:final-aci}

La tabla~\ref{tab:final-aci} compara el intervalo propio, la ventana de
residuos sin adaptación ($\gamma = 0$) y ACI con $\gamma = 0{,}001$.

\begin{table}[htbp]
	\centering
	\caption{Intervalo propio, ventana sin adaptación y ACI ($\gamma = 0{,}001$) en 2024 (8784 horas)}
	\label{tab:final-aci}
	\small
	\setlength{\tabcolsep}{3.5pt}
	\begin{tabular}{lrrrr}
		\hline
		 & \textbf{Ingenuo} & \textbf{SARIMAX} & \textbf{LightGBM} & \textbf{N-HiTS} \\
		\hline
		\textit{Cobertura (\%)} & & & & \\
		\quad Propio & 92,3 & 93,3 & 76,1 & 78,7 \\
		\quad Ventana ($\gamma = 0$) & 86,9 & 87,3 & 88,2 & 89,4 \\
		\quad ACI & 88,0 & 88,6 & 89,1 & 89,4 \\
		\quad IC del 95\,\% de ACI & 85,4--90,6 & 86,1--91,0 & 87,4--90,7 & 87,7--91,1 \\
		\textit{Anchura media (MW)} & & & & \\
		\quad Propio & 5948 & 5717 & 2221 & 2081 \\
		\quad Ventana ($\gamma = 0$) & 5201 & 4796 & 2359 & 2823 \\
		\quad ACI & 5594 & 5175 & 2456 & 2845 \\
		\textit{Puntuación de intervalo (MW)} & & & & \\
		\quad Propio & 8811 & 7392 & 4418 & 4157 \\
		\quad Ventana ($\gamma = 0$) & 9048 & 8115 & 3300 & 3751 \\
		\quad ACI & 9140 & 8144 & 3314 & 3768 \\
		\textit{ACI $-$ propio (MW)} & & & & \\
		\quad Puntuación de intervalo & $+$329 & $+$752 & $-$1104 & $-$388 \\
		\quad Lectura & {\footnotesize no distinguible} & {\footnotesize apreciable} & {\footnotesize apreciable} & {\footnotesize apreciable} \\
		\quad CRPS$_{19}$ & $+$14 & $-$1 & $-$21 & $-$16 \\
		\quad Lectura & {\footnotesize no distinguible} & {\footnotesize no distinguible} & {\footnotesize apreciable} & {\footnotesize apreciable} \\
		\textit{ACI $-$ ventana (MW)} & & & & \\
		\quad Puntuación de intervalo & $+$92 & $+$29 & $+$15 & $+$18 \\
		\quad Lectura & {\footnotesize no distinguible} & {\footnotesize no distinguible} & {\footnotesize no distinguible} & {\footnotesize apreciable} \\
		\textit{Fiabilidad de ACI} & & & & \\
		\quad Desviación máxima & 0,030 & 0,016 & 0,009 & 0,006 \\
		\hline
	\end{tabular}
\end{table}


Para LightGBM y N-HiTS, ACI lleva la cobertura del 76,1 y del 78,7\,\% al 89,1
y al 89,4\,\%, y baja la puntuación de intervalo un 25,0\,\% y un 9,3\,\%
respecto a la del intervalo propio, con diferencias apreciables. Pero la ventana
sin adaptación ya consigue casi lo mismo (3300 y 3751~MW frente a 3314 y
3768~MW con ACI): la mejora viene de calcular el intervalo con los residuos de
las últimas ocho semanas. La adaptación de $\alpha_t$ no aporta nada: la
diferencia de puntuación de intervalo entre ACI y la ventana es de $+15$ y
$+18$~MW, no distinguible en LightGBM y apreciable en N-HiTS pero a favor de la
ventana (un 0,5\,\%). Es el mismo resultado negativo que en 2023.

En los dos modelos clásicos, la ventana queda algo por debajo de la banda
(86,9\,\% en el ingenuo y 87,3\,\% en el SARIMAX) y ACI los sube al 88,0 y al
88,6\,\%, dentro de la banda; por la regla de lectura, ACI mejora ahí la
calibración respecto a la ventana, aunque con intervalos más anchos y sin mejora
de la puntuación de intervalo (diferencias no distinguibles de $+92$ y $+29$~MW).
Frente al intervalo propio, ACI empeora la puntuación del SARIMAX de forma
apreciable ($+752$~MW) y no cambia la del ingenuo (diferencia no distinguible de
$+329$~MW): el intervalo que ya sobrecubría es mejor, en puntuación, que el
construido con ocho semanas de residuos. Con ACI, la desviación máxima de
fiabilidad de los cuatro modelos es de 0,030 o menos.

La tabla~\ref{tab:final-gamma} muestra la sensibilidad a $\gamma$, solo
descriptiva. Con $\gamma = 0{,}001$ la puntuación de intervalo es entre un 0,4 y un
1,0\,\% peor que con la ventana en los cuatro modelos y no hay horas con límite
infinito. Con $\gamma = 0{,}02$ la puntuación baja respecto a la ventana en el
ingenuo (0,950) y en el SARIMAX (0,991), pero sube en los dos de aprendizaje
automático, y hay 396 y 408 horas con límite infinito en los dos clásicos
(el 4,5 y el 4,6\,\% de las horas) y 210 y 156 en los otros dos. Por la regla de
lectura, más del 2\,\% de horas infinitas es un defecto del método, que aquí
se da en tres de los cuatro modelos. La puntuación de la regla de elección de
$\gamma$, recalculada con 2024, es 1,0057, 1,0275, 1,0034 y 0,9965 para
$\gamma$ = 0,001, 0,005, 0,01 y 0,02. Aunque el último valor es menor que uno,
$\gamma$ no se reelige.

\begin{table}[htbp]
	\centering
	\caption{Sensibilidad a $\gamma$ en 2024: puntuación de intervalo relativa a la ventana ($\gamma = 0$) y horas con límite infinito}
	\label{tab:final-gamma}
	\small
	\begin{tabular}{lrrrrrrrr}
		\hline
		 & \multicolumn{4}{c}{\textbf{IS($\gamma$) / IS(0)}} & \multicolumn{4}{c}{\textbf{Horas infinitas}} \\
		$\gamma$ & Ing. & SAR. & LGBM & N-HiTS & Ing. & SAR. & LGBM & N-HiTS \\
		\hline
		0,001 & 1,010 & 1,004 & 1,005 & 1,005 & 0 & 0 & 0 & 0 \\
		0,005 & 1,027 & 1,063 & 1,015 & 1,004 & 114 & 60 & 0 & 0 \\
		0,010 & 0,979 & 0,998 & 1,025 & 1,011 & 186 & 198 & 60 & 30 \\
		0,020 & 0,950 & 0,991 & 1,024 & 1,021 & 396 & 408 & 210 & 156 \\
		\hline
	\end{tabular}
\end{table}


\subsection{Síntesis}\label{sec:final-cal-sintesis}

La calibración de 2024 deja cuatro resultados. Primero, los intervalos propios de
LightGBM y de N-HiTS siguen sin estar calibrados (76,1 y 78,7\,\% para un 90\,\%
nominal), algo menos que en 2023, y los de los dos modelos clásicos sobrecubren
por poco. Segundo, calcular el intervalo con los residuos de las últimas ocho
semanas lleva a los dos modelos de aprendizaje automático a cerca del 89\,\% y
reduce su puntuación de intervalo un 25 y un 9\,\%. Tercero, la adaptación de ACI
no aporta nada sobre esa ventana: el resultado negativo de 2023 se confirma.
Cuarto, en el ingenuo y el SARIMAX ni la ventana ni ACI mejoran, en valor nominal,
la puntuación del intervalo propio, y ACI la empeora de forma apreciable en el
SARIMAX; solo acercan algo la cobertura a la banda.

Entre los intervalos construidos con ACI, los cuatro quedan dentro de la banda
del 88 al 92\,\%, y el de LightGBM es el más estrecho (2456~MW de anchura media
frente a 2845 en N-HiTS), con una puntuación de intervalo nominalmente menor
(3314 frente a 3768~MW). Esa última comparación no se ha contrastado. Los
resultados de calibración tienen las mismas limitaciones que los de 2023: los
residuos iniciales de ACI son los de los modelos de 2023, y el $\gamma$ se eligió
con 2023.

% =====================================================================
% Sección 5.6 redactada el 06/11/2026 (sesión R4). Cifras de
% data/resultados/final/desglose/ (src/30_desglose_oe8_2024.py; protocolo,
% apartado 16.14). Las tablas se generan desde los CSV.
% =====================================================================
\section{Desglose por hora, tipo de día y régimen}\label{sec:desglose}

Esta sección responde al objetivo~OE8: desglosar la precisión y la calibración
de 2024 por hora del día y por tipo de día, e identificar los regímenes en los
que cambia el modelo preferible. El desglose no entrena ni reajusta nada: lee
las previsiones, los intervalos y los resultados de ACI de 2024 que ya se
habían guardado y los agrupa. Se calculó con un único script, una sola vez, con
un candado que impide repetirlo, y ninguno de sus resultados cambia un modelo,
un parámetro ni el $\gamma$ de ACI. Lo que sigue es descripción, contraste con
reglas fijadas de antemano o hipótesis sin comprobar.

La sección empieza con el diseño del desglose, sigue con la precisión por
régimen, con el modelo preferible en cada uno y con la calibración, y termina
con una síntesis.

% ---------------------------------------------------------------------
\subsection{Diseño del desglose}\label{sec:desg-diseno}

\paragraph{Regímenes.} Los 366 días de 2024 se reparten, por el calendario de
la sección~\ref{sec:datos-calendario}, en 253 laborables, 102 días de fin de
semana, 10 festivos y 1 puente. Los festivos son el 1 y el 6 de enero (este,
sábado), el 28 y el 29 de marzo, el 1 de mayo, el 15 de agosto, el 12 de octubre
(sábado), el 1 de noviembre, el 6 de diciembre y el 25 de diciembre; un festivo
en sábado cuenta como festivo. El puente es el viernes 16 de agosto. El 8 de
diciembre, domingo y fiesta nacional, figura como domingo corriente por la
limitación del calendario declarada en la
sección~\ref{sec:datos-limitaciones}. Además de los tipos de día, se separa el
periodo de Navidad (del 23 de diciembre al 6 de enero), que en 2024 suma 15 días,
y se analizan los cuatro bloques horarios B1 a B4 de la
sección~\ref{sec:prob-principio}, el cruce de los tres tipos de día por los
cuatro bloques y el día de la semana.

\paragraph{Qué se conocía ya de 2024.} Al presentar las secciones~\ref{sec:evaluacion-final}
y~\ref{sec:calibracion-final} ya se habían leído el error por bloque y por mes y
la cobertura por bloque y por mes. Eso se repite aquí como control de
reproducción, y las cifras por bloque coinciden con las de aquellas secciones.
Lo que no se había visto es el desglose por hora local, por tipo de día, por
cruce, por día de la semana y por periodo navideño.

\paragraph{Reglas de lectura.} Cada hora local tiene 366 observaciones, de
modo que el intervalo de confianza de una cobertura por hora es de unos tres
puntos y las 24 horas se tratan solo como descripción. Con el puente, un solo
día, no se calcula ningún intervalo ni contraste. Para el resto se contrasta una
celda cada vez. Si el modelo de menor MAE de la celda es LightGBM, el preferible
global, se compara con el segundo; si es otro, se compara LightGBM con ese. La
diferencia de MAE se estima por \textit{bootstrap} pareado de días (1000
réplicas, semilla 42), con ajuste de Holm dentro de cada familia de contrastes
(4 bloques, 3 tipos de día y 12 cruces), y se lee como apreciable si el
$p$-valor ajustado es inferior a 0,05. \textbf{Hay cambio de régimen} en una
celda solo si el mejor modelo es distinto de LightGBM y el intervalo de la
diferencia excluye el cero tras el ajuste; si el mejor es otro pero el
intervalo incluye el cero, se dice que no se distingue y no se cuenta como
cambio. Un cambio solo se afirma en regímenes de al menos 10 días. Para el
sesgo, la calibración y la asimetría de los fallos valen las reglas de las
secciones~\ref{sec:contraste-dm} y~\ref{sec:prob-lectura}: un sesgo es apreciable
por encima de 150~MW, una cobertura está bien calibrada entre el 88 y el
92\,\%, y la asimetría de los fallos es marcada si su cociente supera 2 o baja
de 0,5. Una diferencia de cobertura entre dos regímenes se lee como apreciable
solo si supera cinco puntos y los intervalos de confianza no se solapan.

\paragraph{Controles.} Antes de dar por buenas las salidas, el script comprobó
los recuentos de días, la reproducción del MAE y de las coberturas por bloque
de la evaluación final y que los MAE por hora, por tipo de día y por cruce
reproducen el MAE global. Comprobó también el contraste con una diferencia
sintética: sin diferencia, el $p$-valor fue inferior a 0,05 en el 3\,\% de 60
repeticiones, y con una diferencia de $+60$~MW, el $p$-valor fue de 0,002.

\paragraph{Riesgo de azar.} Con 24 horas, 12 cruces y varios regímenes, algún
patrón aparecerá por casualidad. Por eso las horas son solo descripción, los
contrastes llevan ajuste de Holm por familia y un resultado aislado se lee como
hipótesis.

% ---------------------------------------------------------------------
\subsection{Precisión por régimen}\label{sec:desg-precision}

La tabla~\ref{tab:desg-precision} recoge el MAE y el sesgo de cada modelo en
cada régimen. El sesgo es real menos previsto, de modo que un signo negativo
indica que se prevé más demanda de la que hubo.

\begin{table}[htbp]
	\centering
	\caption{MAE y sesgo (entre paréntesis) por régimen en 2024 (MW)}
	\label{tab:desg-precision}
	\small
	\begin{tabular}{lrrrrr}
		\hline
		\textbf{Régimen} & \textbf{Días} & \textbf{Ingenuo} & \textbf{SARIMAX} & \textbf{LightGBM} & \textbf{N-HiTS} \\
		\hline
		B1 (00--05 h) & 366 & 745 ($-$1) & 825 ($+$111) & 375 ($+$213) & 537 ($-$246) \\
		B2 (06--11 h) & 366 & 1265 ($-$6) & 1118 ($+$19) & 578 ($+$193) & 678 ($+$49) \\
		B3 (12--17 h) & 366 & 1649 ($-$15) & 1304 ($-$158) & 792 ($-$84) & 762 ($-$223) \\
		B4 (18--23 h) & 366 & 1212 ($-$5) & 1112 ($+$52) & 549 ($+$218) & 685 ($+$255) \\
		\hline
		Laborable & 253 & 1218 ($+$140) & 1074 ($+$94) & 582 ($+$183) & 677 ($+$8) \\
		Fin de semana & 102 & 944 ($+$25) & 962 ($+$40) & 525 ($+$57) & 638 ($-$144) \\
		Festivo & 10 & 3672 ($-$3589) & 2813 ($-$2567) & 796 ($-$155) & 655 ($-$255) \\
		Puente (16/8) & 1 & 4639 ($-$4639) & 934 ($+$6) & 1079 ($-$1079) & 624 ($+$200) \\
		\hline
		Navidad (23/12--6/1) & 15 & 2374 ($-$1700) & 1822 ($-$303) & 597 ($-$208) & 762 ($-$355) \\
		Fuera de Navidad & 351 & 1168 ($+$66) & 1058 ($+$19) & 573 ($+$150) & 661 ($-$28) \\
		\hline
	\end{tabular}
\end{table}

\paragraph{Por hora.} La figura~\ref{fig:desg-hora} muestra el MAE de cada modelo
por hora local del día objetivo. Es una descripción, sin contrastes. El máximo
del error está a las 13 h en el ingenuo (1689~MW), a las 15 h en el SARIMAX
(1333~MW), a las 14 h en LightGBM (849~MW) y a las 12 h en N-HiTS (866~MW), que
tiene un segundo máximo a las 21 h (850~MW). LightGBM tiene el menor MAE en 19
de las 24 horas, y N-HiTS en cinco, las de las 13 a las 17 h. Con un único
origen diario, la hora del día y la anticipación van unidas
(sección~\ref{sec:val-perfil}), y la figura no permite separarlas.

\begin{figure}[htbp]
	\centering
	\includegraphics[width=0.85\textwidth]{desglose_2024_mae_hora.png}
	\caption{MAE de los cuatro modelos en 2024 por hora local del día objetivo.}
	\label{fig:desg-hora}
\end{figure}

\paragraph{Por bloque y por tipo de día.} Entre laborables y fines de semana las
diferencias son moderadas: el MAE del ingenuo baja de 1218 a 944~MW, el de
LightGBM de 582 a 525~MW y el de N-HiTS de 677 a 638~MW, mientras que el del
SARIMAX baja menos (de 1074 a 962~MW). En los festivos, en cambio, los modelos
se separan. El ingenuo falla 3672~MW de media, unas tres veces su error de
laborable, y el SARIMAX 2813~MW, unas 2,6 veces el suyo; los dos tienen un
sesgo de $-3589$ y $-2567$~MW, es decir, prevén un día laborable cuando la
demanda es la de un festivo, como en 2023
(sección~\ref{sec:val-dias}). LightGBM falla 796~MW, un 37\,\% más que en un
laborable, y N-HiTS 655~MW, que no es más que en un laborable (677~MW). El sesgo
de festivos es negativo en los cuatro modelos, y en LightGBM y N-HiTS
($-155$ y $-255$~MW) supera el umbral de 150~MW. El cruce más
difícil es el de los festivos por la franja de la mañana (B2), con 5203, 4310,
864 y 739~MW en el ingenuo, el SARIMAX, LightGBM y N-HiTS.

\paragraph{Navidad.} El MAE dentro del periodo de Navidad (15 días) es 2,03 veces
el de fuera en el ingenuo y 1,72 veces en el SARIMAX, frente a 1,04 en LightGBM y
1,15 en N-HiTS. El sesgo dentro de Navidad es negativo en los cuatro modelos.
Es una descripción: son 15 días de un solo año.

\paragraph{Día de la semana, puente y domingo festivo.} Por día de la semana, el
MAE de LightGBM va de 519~MW los domingos a 617~MW los miércoles, y el de N-HiTS,
de 599~MW los viernes a 791~MW los lunes, con un sesgo de $+177$~MW ese día. El
puente del 16 de agosto es un solo día y no se generaliza: el ingenuo falla
4639~MW, con sesgo negativo porque copia el viernes anterior, laborable;
el SARIMAX, 934~MW; LightGBM, 1079~MW, y N-HiTS, 624~MW. El domingo festivo del
8 de diciembre, tratado como domingo corriente, tiene un MAE de 496, 745, 527 y
449~MW en el ingenuo, el SARIMAX, LightGBM y N-HiTS, frente a 944, 962, 525 y
638~MW de media en los fines de semana: la omisión del calendario no se nota en
ese día.

\paragraph{Peores días.} Los 10 peores días de cada modelo, por MAE diario,
están cerca de festivos y puentes. Un cálculo aparte sobre las fechas, que no
está en el script y es solo descriptivo, cuenta cuántos están a siete días o
menos de un festivo o del puente: los 10 del ingenuo, los 10 del SARIMAX, 8 de
los de LightGBM y 4 de los de N-HiTS. Como el 32,5\,\% de los días de 2024 (119
de 366) cumple esa condición, la concentración es muy marcada en los dos modelos
clásicos, menos en LightGBM, aunque 8 de 10 supera el 32,5\,\%, y no se aprecia en N-HiTS (4 de 10). El peor día del ingenuo
es el 25 de diciembre (9212~MW) y el del SARIMAX, el 8 de enero (7199~MW); los
de LightGBM y N-HiTS son el 12 de octubre (1692~MW), un festivo en sábado, y el
10 de marzo (2306~MW), un domingo.

% ---------------------------------------------------------------------
\subsection{Modelo preferible y cambio de régimen}\label{sec:desg-regimen}

La tabla~\ref{tab:desg-pref} reúne los 19 contrastes de la sección anterior. En
cada fila, la diferencia es el MAE de LightGBM menos el del modelo con que se
compara, de modo que un valor negativo favorece a LightGBM.

\begin{table}[htbp]
	\centering
	\caption{Modelo preferible por celda y contraste con LightGBM (MW; $p$ ajustado por Holm dentro de cada familia)}
	\label{tab:desg-pref}
	\small
	\setlength{\tabcolsep}{3pt}
	\begin{tabular}{lrrrr>{\raggedright\arraybackslash}p{4.1cm}}
		\hline
		\textbf{Celda} & \textbf{Días} & \textbf{Diferencia} & \textbf{IC del 95\,\%} & \textbf{$p$ ajustado} & \textbf{Lectura} \\
		\hline
		B1 (00--05 h) & 366 & $-$162 & $-$211 a $-$114 & 0,008 & LightGBM, apreciable \\
		B2 (06--11 h) & 366 & $-$100 & $-$143 a $-$57 & 0,008 & LightGBM, apreciable \\
		B3 (12--17 h) & 366 & $+$30 & $-$21 a $+$77 & 0,242 & N-HiTS nominal, no distinguible \\
		B4 (18--23 h) & 366 & $-$136 & $-$181 a $-$92 & 0,008 & LightGBM, apreciable \\
		\hline
		Laborable & 253 & $-$95 & $-$139 a $-$52 & 0,006 & LightGBM, apreciable \\
		Fin de semana & 102 & $-$113 & $-$183 a $-$53 & 0,006 & LightGBM, apreciable \\
		Festivo & 10 & $+$141 & $-$57 a $+$403 & 0,222 & N-HiTS nominal, no distinguible \\
		\hline
		Laborable $\times$ B1 & 253 & $-$165 & $-$226 a $-$102 & 0,024 & LightGBM, apreciable \\
		Laborable $\times$ B2 & 253 & $-$119 & $-$178 a $-$56 & 0,024 & LightGBM, apreciable \\
		Laborable $\times$ B3 & 253 & $+$30 & $-$27 a $+$81 & 0,893 & N-HiTS nominal, no distinguible \\
		Laborable $\times$ B4 & 253 & $-$126 & $-$180 a $-$69 & 0,024 & LightGBM, apreciable \\
		Fin de semana $\times$ B1 & 102 & $-$154 & $-$265 a $-$61 & 0,024 & LightGBM, apreciable \\
		Fin de semana $\times$ B2 & 102 & $-$79 & $-$178 a $+$4 & 0,408 & LightGBM, no distinguible \\
		Fin de semana $\times$ B3 & 102 & $-$25 & $-$118 a $+$64 & 1,000 & LightGBM, no distinguible \\
		Fin de semana $\times$ B4 & 102 & $-$193 & $-$275 a $-$114 & 0,024 & LightGBM, apreciable \\
		Festivo $\times$ B1 & 10 & $-$185 & $-$357 a $-$16 & 0,224 & LightGBM, no distinguible \\
		Festivo $\times$ B2 & 10 & $+$126 & $-$187 a $+$462 & 1,000 & N-HiTS nominal, no distinguible \\
		Festivo $\times$ B3 & 10 & $+$384 & $-$56 a $+$886 & 0,408 & N-HiTS nominal, no distinguible \\
		Festivo $\times$ B4 & 10 & $+$238 & $-$35 a $+$480 & 0,408 & N-HiTS nominal, no distinguible \\
		\hline
	\end{tabular}
\end{table}

\textbf{No hay ningún cambio de régimen distinguible.} LightGBM se mantiene como
preferible de forma apreciable en 10 de las 19 celdas: tres de los cuatro
bloques (B1, B2 y B4), los laborables, los fines de semana, los cruces de
laborable con B1, B2 y B4 y los de fin de semana con B1 y B4. Se mantiene como
mejor sin distinguirse del segundo en otras tres (fin de semana por B2 y por
B3, y festivo por B1). En las seis restantes el menor MAE nominal es de N-HiTS,
pero la diferencia no se distingue de cero: la franja de la tarde (B3), con una
diferencia de $+30$~MW, el laborable por B3, los festivos ($+141$~MW, con un
intervalo de $-57$ a $+403$~MW) y los festivos por B2, B3 y B4. El puente, con un
solo día, queda como descripción.

Hay que leer ese «no se distingue» con prudencia. Con 10 días, el contraste de
los festivos tiene poca potencia, de modo que la falta de una diferencia
distinguible no equivale a que los dos modelos sean igual de buenos. Tampoco se
distingue en la tarde de los laborables (253 días), pero ahí la diferencia
nominal es de solo 30~MW.

% ---------------------------------------------------------------------
\subsection{Calibración por régimen}\label{sec:desg-cal}

La tabla~\ref{tab:desg-cob} recoge la cobertura del intervalo del 90\,\% por
tipo de día, con el intervalo propio y con ACI.

\begin{table}[htbp]
	\centering
	\caption{Cobertura del intervalo del 90\,\% por tipo de día en 2024 (\%): intervalo propio / ACI ($\gamma = 0{,}001$)}
	\label{tab:desg-cob}
	\small
	\begin{tabular}{lrrrr}
		\hline
		\textbf{Tipo de día} & \textbf{Ingenuo} & \textbf{SARIMAX} & \textbf{LightGBM} & \textbf{N-HiTS} \\
		\hline
		Laborable & 92,4 / 87,4 & 93,7 / 89,3 & 76,8 / 88,5 & 78,3 / 88,4 \\
		Fin de semana & 97,4 / 94,8 & 94,8 / 90,8 & 76,1 / 92,5 & 78,9 / 91,9 \\
		Festivo & 47,5 / 42,5 & 65,8 / 50,0 & 65,4 / 74,2 & 86,2 / 90,8 \\
		Puente (24 horas) & 0,0 / 0,0 & 100,0 / 83,3 & 16,7 / 54,2 & 91,7 / 83,3 \\
		\hline
	\end{tabular}
\end{table}

En \textbf{laborables y fines de semana} los resultados se parecen a los del año
completo (sección~\ref{sec:calibracion-final}). Los intervalos propios de
LightGBM y N-HiTS infracubren en ambos (entre el 76,1 y el 78,9\,\%), los del
ingenuo y el SARIMAX sobrecubren y ACI deja a los cuatro modelos cerca de la
banda del 88 al 92\,\%, con tres excepciones: el ingenuo con ACI queda por
debajo en laborables (87,4\,\%) y por encima en fines de semana (94,8\,\%), y
LightGBM con ACI queda apenas por encima en fines de semana (92,5\,\%). Esa diferencia de 7,4
puntos entre los dos tipos de día, con intervalos de confianza que no se
solapan, es la única diferencia de calibración entre regímenes que se lee como
apreciable: con el intervalo propio, el ingenuo cubre un 92,4\,\% en laborables
y un 97,4\,\% en fines de semana, que son 4,99 puntos y, por la regla, no
alcanzan el umbral.

En los \textbf{festivos} la calibración se rompe en los modelos clásicos. El
intervalo propio del ingenuo cubre el 47,5\,\% de las horas y el del SARIMAX el
65,8\,\%, y ACI no lo repara: el ingenuo queda en el 42,5\,\% y el SARIMAX en el
50,0\,\%. Con ACI, el intervalo de confianza de LightGBM (60,0 a 89,2\,\%)
tampoco incluye el 90\,\%; solo el de N-HiTS lo incluye (80,8 a 98,3\,\%), con
una cobertura del 90,8\,\%. Con 10 días los intervalos de confianza son muy
anchos y la comparación es descriptiva. La puntuación de intervalo y el
CRPS$_{19}$ de los festivos, con el intervalo propio, ordenan los modelos como
el error: 35\,733 y 2984~MW en el ingenuo, 20\,906 y 2105 en el SARIMAX, 9601 y
658 en LightGBM, y 3397 y 528 en N-HiTS; con ACI, 5542 y 615 en LightGBM y 3573
y 481 en N-HiTS. En el cruce de festivos con la mañana (B2), la cobertura
propia es del 31,7, 33,3, 53,3 y 76,7\,\% en los cuatro modelos.

Una causa posible de que ACI no repare los festivos de los modelos clásicos es
que su retroalimentación llega con 24 a 48 horas de retraso y que los festivos
son días aislados, de modo que los errores del día anterior no avisan del
siguiente. Es coherente con los datos, pero no se ha comprobado. En el puente,
un solo día, las coberturas son del 0, el 100, el 16,7 y el 91,7\,\% (propio),
sin intervalos. No se ha analizado la cobertura por hora local: queda en el
repositorio, y su lectura sería posterior y descriptiva.

% ---------------------------------------------------------------------
\subsection{Síntesis}\label{sec:desg-sintesis}

El desglose deja cuatro resultados. Primero, no se encuentra ningún régimen en
el que cambie de forma distinguible el modelo preferible: LightGBM se mantiene
como mejor, de forma apreciable, en los regímenes con datos suficientes. Segundo,
el efecto fuerte del régimen no está entre LightGBM y N-HiTS, sino en los
modelos clásicos, que en los festivos multiplican por unas tres veces su error
de un laborable, quedan lejos de la cobertura del 90\,\% y no la recuperan con
ACI. Tercero, \textbf{no se concluye que N-HiTS sea mejor en épocas difíciles}:
su ventaja nominal en festivos y en la franja de la tarde no se distingue de cero,
y LightGBM tiene un error menor en Navidad (597 frente a 762~MW). Queda como
hipótesis, con 10 festivos de un solo año, que N-HiTS sea más robusto ante los
festivos, porque es el único modelo con una cobertura compatible con el 90\,\%
en ellos y tiene la menor puntuación de intervalo, sin contraste. Cuarto, las
causas de todo lo anterior (el fallo de los clásicos en festivos, la ventaja
nominal de N-HiTS por la tarde, la falta de efecto de ACI) son hipótesis sin
comprobar, y las conclusiones describen un solo año.

% =====================================================================
% Sección 5.7 redactada el 06/11/2026 (sesión R4). No introduce cifras
% nuevas: reúne las de 5.2 a 5.6. Las explicaciones de los resultados se
% declaran como hipótesis cuando no se han comprobado.
% =====================================================================
\section{Discusión}\label{sec:discusion}

Esta sección interpreta los resultados de las anteriores. Responde primero a la
pregunta de investigación de la sección~\ref{sec:aportacion}, separa después lo
que los datos muestran de lo que solo se conjetura sobre sus causas, recoge los
resultados negativos y termina con los límites que acotan lo que se puede
afirmar. No se presentan cifras nuevas: todas proceden de las secciones
anteriores, a las que se remite.

% ---------------------------------------------------------------------
\subsection{Respuesta a la pregunta de investigación}\label{sec:disc-respuesta}

La pregunta tiene tres partes. La primera es si los modelos de aprendizaje
automático y profundo mejoran, de forma estadísticamente significativa, al
SARIMAX con estacionalidad múltiple en la previsión del día siguiente. En 2024
\textbf{sí lo hacen}, con las condiciones del protocolo: el MAE de LightGBM es
de 573,6~MW y el de N-HiTS de 665,6~MW, frente a los 1089,7~MW del SARIMAX, es
decir, un 47 y un 39\,\% menos. Los dos contrastes de Diebold--Mariano tienen un
p-valor ajustado inferior a 0,001 y se mantienen con la pérdida cuadrática y con 3
y 14 retardos (sección~\ref{sec:final-dm}). Entre los dos modelos de aprendizaje
automático, LightGBM es apreciablemente mejor que N-HiTS, por 92~MW de pérdida
media diaria. Hay un matiz que reduce el alcance de esta respuesta: el SARIMAX
no se distingue del ingenuo estacional con la regla fijada de antemano. Superar
al SARIMAX es, por tanto, superar a un comparador que no mejora de forma
demostrable a la referencia más simple, y la comparación más informativa es la
de los dos modelos de aprendizaje automático con el ingenuo (diferencias de 644
y 552~MW de pérdida media diaria).

La segunda parte pregunta en qué horas y tipos de día lo hacen. La respuesta es
que \textbf{en casi todos}. LightGBM tiene el menor MAE en 19 de las 24 horas
y N-HiTS en cinco, las de las 13 a las 17 h; ningún modelo clásico es el
mejor en ninguna hora. En ninguno de los 19 contrastes del desglose cambia de
forma distinguible el modelo preferible, de modo que no se encuentra un régimen
en el que un modelo clásico sea la mejor opción
(sección~\ref{sec:desg-regimen}). Donde la distancia entre familias es mayor es
en los festivos: el MAE de los modelos clásicos es de 3672 y 2813~MW y el de los
de aprendizaje automático, de 796 y 655~MW. Esta es la respuesta más clara a «en
qué situaciones», pero se apoya en 10 festivos de un solo año
(sección~\ref{sec:desg-precision}).

La tercera parte pregunta qué familia proporciona intervalos mejor calibrados, y
la respuesta es menos sencilla. Con los intervalos propios de cada modelo
ninguna familia queda en la banda del 88 al 92\,\%, pero las dos fallan en
sentidos distintos. Los intervalos de los modelos clásicos cubren algo de más
(92,3 y 93,3\,\%), mientras que los cuantiles de LightGBM y N-HiTS cubren
claramente de menos (76,1 y 78,7\,\%). En ese sentido, los modelos clásicos
están mejor calibrados, y lo están a costa de intervalos más anchos. Pero la
calibración no es una propiedad fija de la familia: cuando los intervalos de los
cuatro se construyen con el mismo método, la ventana de residuos recientes o
ACI, los cuatro quedan dentro de la banda
(sección~\ref{sec:final-cal-sintesis}). Lo que se puede decir con los datos es,
por tanto, que \textbf{la calibración depende más de cómo se construye el
intervalo que del modelo}, y que los cuantiles entrenados con la pérdida de
pinball, sin corregir, no son fiables en estos dos modelos. Que LightGBM y
N-HiTS tengan una puntuación de intervalo menor no contradice esto, porque se
debe en parte a que sus intervalos son estrechos (sección~\ref{sec:cal-propios}).

Estas respuestas son del problema y el protocolo estudiados. La literatura
revisada es mixta sobre el rendimiento relativo de los métodos estadísticos y de
aprendizaje automático (sección~\ref{sec:evidencia}), y un único sistema, un
único año de evaluación y un modelo por familia no la resuelven. El resultado
concuerda con la ventaja del aprendizaje automático en este problema concreto,
pero no permite extenderla a otros.

% ---------------------------------------------------------------------
\subsection{Qué se sabe y qué se conjetura}\label{sec:disc-hipotesis}

Los resultados describen qué ocurre; las explicaciones de por qué ocurre, en
general, no se han comprobado. Se separan aquí las dos cosas. Las que siguen son
hipótesis, y ninguna se ha contrastado.

\paragraph{Exceso de previsión al mediodía y autoconsumo.} En 2023 los tres
modelos ajustados prevén entre las 12:00 y las 18:00 más demanda de la que hubo
($-508$, $-411$ y $-313$~MW de sesgo medio en el SARIMAX, LightGBM y N-HiTS), y en
LightGBM el exceso es mayor de abril a septiembre ($-507$~MW) que de noviembre a
febrero ($-210$~MW) (sección~\ref{sec:val-perfil}). Es compatible con el
autoconsumo fotovoltaico, que reduce la demanda observada en las horas de sol y
no está entre las variables de los modelos, que solo conocen la temperatura y el
calendario. Pero la misma pauta podría deberse a otra variable estacional no
incluida, y no se ha medido el autoconsumo. En 2024 el sesgo de esa franja es
menor ($-158$, $-84$ y $-223$~MW, tabla~\ref{tab:desg-precision}), lo que
coincide con el reajuste mensual, que podría absorber parte del crecimiento del
autoconsumo, pero también con el cambio de año, y el diseño no permite separar
ambas cosas.

\paragraph{El perfil aplanado de N-HiTS.} En 2023 la amplitud del perfil diario
de N-HiTS es el 88\,\% de la real, frente al 101--102\,\% de los otros dos modelos
ajustados (sección~\ref{sec:val-perfil}). Una causa posible está en la
arquitectura, que genera la previsión interpolando desde representaciones de
baja resolución \parencite{challu2023nhits}, lo que favorece perfiles suaves; no
se ha comprobado. Además, el análisis del perfil no se ha repetido en 2024.

\paragraph{Cuantiles estrechos.} Los cuantiles de LightGBM y N-HiTS se ajustan
con los datos de entrenamiento, donde el error del modelo suele ser menor que
fuera de la muestra, y esto produciría intervalos más estrechos de lo debido. Es
coherente con que la ventana de residuos recientes, que son errores fuera de
muestra, los corrija. No se ha comparado la cobertura dentro y fuera de la
muestra de entrenamiento, que es lo que lo comprobaría.

\paragraph{Anchura de los intervalos del SARIMAX.} Su intervalo gaussiano
homocedástico tiene la misma anchura en todas las horas, lo que se sigue de su
construcción. Esa anchura no refleja que el error depende de la hora, y en
2023 resultó demasiado ancha de madrugada y justa por la tarde
(sección~\ref{sec:cal-propios}).

\paragraph{Meses y días difíciles.} En 2023, agosto es el mes de menor cobertura del
ingenuo y de LightGBM, y en 2024 es uno de los tres meses de mayor error del
ingenuo. Esto coincide
con el periodo vacacional, en el que la demanda de una semana se parece menos a
la de la anterior, pero esta explicación se propone sin haberse comprobado. Los
modelos clásicos fallan en festivos con un sesgo de $-3589$ y $-2567$~MW, es
decir, prevén más demanda de la que hubo; por qué no los corrige ACI, y si es porque su
retroalimentación llega con 24 a 48 horas de retraso, es una hipótesis.

\paragraph{Robustez de N-HiTS.} Que N-HiTS sea más robusto en festivos y en
épocas difíciles es una hipótesis que sugieren su cobertura del 90,8\,\% y su
menor puntuación de intervalo en ellos, pero no se sostiene con el contraste:
su ventaja de error en festivos no se distingue de cero y LightGBM tiene un
error menor en Navidad (sección~\ref{sec:desg-regimen}).

\paragraph{El empeoramiento del SARIMAX en 2024.} El SARIMAX es el único modelo
cuyo error aumenta entre 2023 y 2024 (un 15,5\,\%), aunque su sesgo se reduce
casi a cero. No se ha identificado la causa. Las explicaciones posibles
(sensibilidad a las características del año, efecto del reajuste mensual, o
que 2024 favorezca más al ingenuo, cuyo error baja un 13\,\%) no se pueden
separar con este diseño.

% ---------------------------------------------------------------------
\subsection{Resultados negativos}\label{sec:disc-negativos}

Cuatro resultados del trabajo son negativos y se mantienen como tales.

El primero es que \textbf{el SARIMAX con estacionalidad múltiple no se distingue
del ingenuo estacional}. Con la regla fijada de antemano (pérdida absoluta y 7
retardos), el p-valor ajustado es 0,087; con la pérdida cuadrática y 3 retardos
sería 0,047. Lo que se informa es la primera lectura, y lo que dicen las
variantes es que, en el mejor caso, el margen es muy pequeño. La especificación
del SARIMAX tiene un papel en esto: en 2023 los mayores cambios los produjeron
la variable de nivel y los retardos de 48 y 168 horas, no el orden del proceso
de los errores (sección~\ref{sec:val-seleccion}), y no se ha explorado si otra
especificación lo acercaría a los modelos de aprendizaje automático. Es una
limitación del comparador y no una propiedad de la familia.

El segundo es que \textbf{los cuantiles de LightGBM y N-HiTS no están
calibrados}, ni en 2023 ni en 2024. Su buen CRPS$_{19}$ no debe leerse como una
buena calibración.

El tercero es que \textbf{ACI no aporta nada sobre la ventana de residuos} en
LightGBM y N-HiTS; en el ingenuo y el SARIMAX ni la ventana ni ACI mejoran la
puntuación del intervalo propio, y ACI la empeora de forma apreciable en el
SARIMAX; en los festivos, tampoco repara la cobertura de los modelos
clásicos. La conclusión es sobre ACI \parencite{gibbs2021adaptive} con el
$\gamma = 0{,}001$ que se eligió con 2023, un valor con el que la adaptación es
casi inerte, y no sobre ACI con otros valores, que no se evaluaron en 2024. Que la
ventana por sí sola baste es un resultado útil: es un método más simple.

El cuarto es que \textbf{el desglose no encuentra un cambio de régimen}: ningún
régimen cambia de forma distinguible el modelo preferible. Para el objetivo OE8
esto es un hallazgo, no una carencia: la ventaja de LightGBM sobre N-HiTS es
estable entre los regímenes con datos suficientes, y la de los dos sobre los
clásicos es grande en todos.

% ---------------------------------------------------------------------
\subsection{Límites de la comparación}\label{sec:disc-limites}

\paragraph{Previsión meteorológica perfecta.} Los modelos usan la temperatura de
reanálisis de la propia hora que prevén, lo que equivale a suponer una previsión
meteorológica sin error (sección~\ref{sec:aportacion}). Los resultados son, por
tanto, una cota superior de lo alcanzable en operación. Con previsiones reales
el error de todos los modelos que usan la temperatura aumentaría, y la ventaja
de los que más dependen de ella podría reducirse. Esa dependencia
no se ha medido en este trabajo. El ingenuo es el único modelo que no usa la temperatura, de modo
que es el que menos se vería afectado, y es posible que su distancia a los
demás se redujera en operación. No se ha cuantificado.

\paragraph{Un sistema, un año y un modelo por familia.} Las conclusiones se
refieren al sistema peninsular, a 2024 y a los cuatro modelos estudiados. Un
solo año tiene un solo patrón de festivos, uno de ellos puente, y los
contrastes sobre 10 festivos tienen intervalos muy anchos. Que LightGBM supere a
N-HiTS no es una conclusión sobre el boosting frente a las redes en general.

\paragraph{Selección con 2023.} Los hiperparámetros de los tres modelos
ajustados y el $\gamma$ de ACI se eligieron con 2023, y esa elección favorece a
los modelos en 2023 pero no en 2024. El efecto es la «maldición del ganador»,
que se declara en la sección~\ref{sec:final-2023}: la comparación de 2023 con
2024 mezcla ese efecto con el del año y el del reajuste mensual, y por eso es
solo descriptiva. No se hizo un control sin reajuste, de modo que no se sabe
cuánto aporta el reajuste mensual a los resultados de 2024.

\paragraph{Contrastes y residuos.} Los contrastes de Diebold--Mariano usan la
pérdida diaria (366 días), que reduce la dependencia entre horas pero no la
elimina, y su resultado para el par SARIMAX--ingenuo depende algo de las
decisiones del contraste. Los residuos iniciales de ACI son los de los modelos
de 2023. El desglose por hora, tipo de día y régimen es descriptivo en lo que no
tiene contraste.

\paragraph{Lo que el trabajo no hace.} No analiza la cobertura por hora local, no
se ha medido el efecto de ninguna de las hipótesis de la sección anterior y no
evalúa otros horizontes ni otros sistemas eléctricos. Quedan como trabajo futuro
en el capítulo~\ref{cap:conclusiones}.
